% =====================================================================
\chapter{Soluciones de los ejercicios}
\label{ap:soluciones}
% =====================================================================

Cada solución se presenta en cinco partes: datos, planteamiento, desarrollo, resultado y comprobación. En los ejercicios de demostración ($\star\star\star$), los pasos del desarrollo, leídos en orden, forman la demostración redactada tal como debería entregarse: cada paso es una frase completa que introduce sus variables y justifica lo que afirma. La comprobación indica una forma independiente de verificar el resultado o señala dónde se ha usado cada hipótesis.

% ---------------------------------------------------------------------
\section{Tema 1}

\begin{solucion}{eje:1-proposiciones}
\datos seis enunciados.
\planteamiento según la \cref{def:proposicion}, un enunciado es proposición si es verdadero o falso, sin depender de variables libres.
\pasos
\begin{enumerate}[label=(\alph*)]
  \item $2^{10}=1024>1000$: proposición \textbf{verdadera}.
  \item Contiene la variable libre $x$: es un predicado, \textbf{no} una proposición (aunque $\forall x\in\R\ (x^{2}\ge 0)$ sí lo sería, y verdadera).
  \item $91=7\cdot 13$: proposición \textbf{falsa}.
  \item Es una orden: \textbf{no} es proposición.
  \item Proposición \textbf{falsa}, porque $x^{2}\ge 0>-1$ para todo real $x$.
  \item Proposición \textbf{verdadera} (\cref{teo:division} con $b=2$).
\end{enumerate}
\resultado proposiciones: (a) V, (c) F, (e) F, (f) V; no son proposiciones (b) y (d).
\comprobacion (b) se convierte en proposición al fijar $x$ o al cuantificarla, como se verá en el \cref{tema:2}.
\end{solucion}

\begin{solucion}{eje:1-traducir}
\datos $P$: $n$ par; $Q$: $3\divides n$; $R$: $6\divides n$.
\planteamiento «$A$ solo si $B$» es $A\Rightarrow B$; «$A$ no es suficiente para $B$» es $\neg(A\Rightarrow B)$.
\pasos
\begin{enumerate}[label=(\alph*)]
  \item $R\Leftrightarrow(P\wedge Q)$. Es verdadera para todo $n$: $6\divides n$ si y solo si $2\divides n$ y $3\divides n$.
  \item $P\Rightarrow R$. Falsa en general: $n=2$ es par y no es múltiplo de~6.
  \item $\neg(P\Rightarrow R)\eqlog P\wedge\neg R$. No es verdadera para todo $n$ (falla con $n=6$); lo que es cierto es que existe un $n$ que la cumple ($n=2$).
  \item $\neg P\wedge\neg Q$, o equivalentemente $\neg(P\vee Q)$ por De Morgan.
\end{enumerate}
\resultado (a) $R\Leftrightarrow(P\wedge Q)$, siempre verdadera; (b) $P\Rightarrow R$, no siempre; (c) $P\wedge\neg R$, no siempre; (d) $\neg P\wedge\neg Q$.
\comprobacion con $n=2$: $P=\V$, $Q=\F$, $R=\F$; (a) da $\F\Leftrightarrow\F=\V$, (b) da $\V\Rightarrow\F=\F$, (c) da $\V$.
\end{solucion}

\begin{solucion}{eje:1-reciproco}
\datos $P$: «$f$ es derivable en $a$»; $Q$: «$f$ es continua en $a$». Se sabe (bachillerato) que $P\Rightarrow Q$ es verdadero.
\planteamiento aplicar la \cref{def:reciproco} y la \cref{prop:leyes}(v).
\pasos
\begin{enumerate}[label=\arabic*.]
  \item Recíproco $Q\Rightarrow P$: «si $f$ es continua en $a$, es derivable en $a$». Falso: $f(x)=\abs{x}$ es continua en $0$ y no es derivable allí (las derivadas laterales valen $-1$ y $1$).
  \item Contrarrecíproco $\neg Q\Rightarrow\neg P$: «si $f$ no es continua en $a$, no es derivable en $a$». Verdadero, por ser equivalente al enunciado original.
  \item Inverso $\neg P\Rightarrow\neg Q$: «si $f$ no es derivable en $a$, no es continua en $a$». Es el contrarrecíproco del recíproco, luego equivalente a él, y también falso con $f(x)=\abs{x}$.
\end{enumerate}
\resultado \textbf{contrarrecíproco verdadero; recíproco e inverso falsos} (contraejemplo $f(x)=\abs{x}$, $a=0$).
\comprobacion el recíproco y el inverso siempre tienen el mismo valor de verdad, pues $Q\Rightarrow P\eqlog\neg P\Rightarrow\neg Q$.
\end{solucion}

\begin{solucion}{eje:1-necesaria}
\datos $T$: «$4\divides n$».
\planteamiento $C$ es suficiente si $C\Rightarrow T$; necesaria si $T\Rightarrow C$. Se busca demostración o contraejemplo para cada implicación.
\pasos
\begin{enumerate}[label=(\alph*)]
  \item $T\Rightarrow$ par: si $n=4k$, $n=2(2k)$. Par $\Rightarrow T$ falla con $n=2$. \textbf{Necesaria, no suficiente.}
  \item $8\divides n\Rightarrow T$: $n=8k=4(2k)$. $T\Rightarrow 8\divides n$ falla con $n=4$. \textbf{Suficiente, no necesaria.}
  \item $T\Rightarrow 16\divides n^{2}$: $n=4k$ da $n^{2}=16k^{2}$. Recíproco: si $n$ es impar, $n^{2}$ es impar; si $n=2m$ con $m$ impar, $n^{2}=4m^{2}$ con $m^{2}$ impar, que no es múltiplo de~16 (sería $4m^{2}=16j$, $m^{2}=4j$ par). Por contrarrecíproco, $16\divides n^{2}\Rightarrow n=2m$ con $m$ par $\Rightarrow 4\divides n$. \textbf{Necesaria y suficiente.}
  \item «Termina en~4» no implica $T$ ($14$) y $T$ no la implica ($8$). \textbf{Ninguna.}
\end{enumerate}
\resultado (a) necesaria; (b) suficiente; (c) necesaria y suficiente; (d) ninguna.
\comprobacion valores: $n=12$ cumple $T$, es par, $144=16\cdot 9$, no termina en~4 y no es múltiplo de~8, coherente con todo lo anterior.
\end{solucion}

\begin{solucion}{eje:1-silogismo}
\datos $\Phi=\bigl((P\Rightarrow Q)\wedge(Q\Rightarrow R)\bigr)\Rightarrow(P\Rightarrow R)$.
\planteamiento construir las $2^{3}=8$ filas (\cref{prop:filas}).
\pasos
\begin{enumerate}[label=\arabic*.]
  \item Tabla, con $H=(P\Rightarrow Q)\wedge(Q\Rightarrow R)$:
  \begin{center}\small
  \begin{tabular}{@{}ccc|ccccc@{}}
  \toprule
  $P$&$Q$&$R$&$P\Rightarrow Q$&$Q\Rightarrow R$&$H$&$P\Rightarrow R$&$\Phi$\\
  \midrule
  V&V&V&V&V&V&V&V\\
  V&V&F&V&F&F&F&V\\
  V&F&V&F&V&F&V&V\\
  V&F&F&F&V&F&F&V\\
  F&V&V&V&V&V&V&V\\
  F&V&F&V&F&F&V&V\\
  F&F&V&V&V&V&V&V\\
  F&F&F&V&V&V&V&V\\
  \bottomrule
  \end{tabular}
  \end{center}
  \item La última columna es siempre $\V$: $\Phi$ es una tautología.
\end{enumerate}
\resultado \textbf{$\Phi$ es una tautología}; por definición, el silogismo hipotético es un razonamiento válido.
\comprobacion sin tabla: $\Phi$ solo podría ser falsa con $P\Rightarrow R$ falsa, es decir, $P=\V$, $R=\F$; entonces $Q\Rightarrow R$ verdadera exige $Q=\F$, y $P\Rightarrow Q$ sería $\V\Rightarrow\F=\F$, luego $H=\F$ y $\Phi=\V$.
\end{solucion}

\begin{solucion}{eje:1-xor}
\datos definición de $\xor$.
\planteamiento comparar columnas de las tablas.
\pasos
\begin{enumerate}[label=\arabic*.]
  \item Tabla (filas $PQ=$ VV, VF, FV, FF):
  \begin{center}\small
  \begin{tabular}{@{}cc|ccccc@{}}
  \toprule
  $P$&$Q$&$P\xor Q$&$P\vee Q$&$\neg(P\wedge Q)$&$(P\vee Q)\wedge\neg(P\wedge Q)$&$\neg(P\Leftrightarrow Q)$\\
  \midrule
  V&V&F&V&F&F&F\\
  V&F&V&V&V&V&V\\
  F&V&V&V&V&V&V\\
  F&F&F&F&V&F&F\\
  \bottomrule
  \end{tabular}
  \end{center}
  \item Las columnas 3, 6 y 7 coinciden fila a fila.
\end{enumerate}
\resultado (a) $P\xor Q$ vale $\F,\V,\V,\F$; (b) y (c) son \textbf{equivalencias ciertas}.
\comprobacion intuitivamente, «exactamente una» es «al menos una» y «no las dos», y también «no tienen el mismo valor».
\end{solucion}

\begin{solucion}{eje:1-equivalencias}
\datos tres pares de fórmulas en $P,Q,R$.
\planteamiento basta estudiar las filas en que alguna fórmula es falsa.
\pasos
\begin{enumerate}[label=(\alph*)]
  \item $P\Rightarrow(Q\Rightarrow R)$ es falsa solo si $P=\V$ y $Q\Rightarrow R$ falsa, es decir, $P=Q=\V$, $R=\F$. $(P\wedge Q)\Rightarrow R$ es falsa solo si $P\wedge Q=\V$ y $R=\F$: la misma fila. \textbf{Equivalentes.}
  \item $P\Rightarrow(Q\Rightarrow R)$ es falsa solo en $(\V,\V,\F)$. $(P\Rightarrow Q)\Rightarrow R$ es falsa cuando $R=\F$ y $P\Rightarrow Q=\V$, es decir, en $(\V,\V,\F)$, $(\F,\V,\F)$ y $(\F,\F,\F)$. \textbf{No equivalentes}: en $P=Q=R=\F$, la primera es $\V$ y la segunda $\F$.
  \item $(P\vee Q)\Rightarrow R$ es falsa si $R=\F$ y $P\vee Q=\V$: filas $(\V,\V,\F)$, $(\V,\F,\F)$, $(\F,\V,\F)$. $(P\Rightarrow R)\wedge(Q\Rightarrow R)$ es falsa si alguna de las dos lo es: $R=\F$ con $P=\V$ o con $Q=\V$, las mismas tres filas. \textbf{Equivalentes.}
\end{enumerate}
\resultado (a) sí; (b) no (asignación $P=Q=R=\F$, también $P=\F,Q=\V,R=\F$); (c) sí.
\comprobacion (a) es la «exportación»: demostrar «si $P$ y $Q$, entonces $R$» equivale a suponer $P$ y luego demostrar «si $Q$, entonces $R$». (c) justifica las demostraciones por casos (\cref{def:casos}).
\end{solucion}

\begin{solucion}{eje:1-algebra}
\datos leyes (i)–(vii) de la \cref{prop:leyes}.
\planteamiento transformar paso a paso indicando la ley usada.
\pasos
\begin{enumerate}[label=\arabic*.]
  \item Distributiva: $P\vee(\neg P\wedge Q)\eqlog(P\vee\neg P)\wedge(P\vee Q)$.
  \item $P\vee\neg P$ es tautología, y $\V\wedge A\eqlog A$: queda $P\vee Q$.
  \item Luego $\neg\bigl(P\vee(\neg P\wedge Q)\bigr)\eqlog\neg(P\vee Q)\eqlog\neg P\wedge\neg Q$ por De Morgan.
  \item Por (iv): $(P\Rightarrow Q)\wedge(P\Rightarrow\neg Q)\eqlog(\neg P\vee Q)\wedge(\neg P\vee\neg Q)$.
  \item Distributiva (sacando $\neg P$): $\eqlog\neg P\vee(Q\wedge\neg Q)$.
  \item $Q\wedge\neg Q$ es contradicción y $A\vee\F\eqlog A$: queda $\neg P$.
\end{enumerate}
\resultado \textbf{ambas equivalencias se cumplen}.
\comprobacion la segunda dice que, si $P$ implicara a la vez $Q$ y $\neg Q$, $P$ sería falsa: es el fundamento de la reducción al absurdo (\cref{tema:4}). Por tabla: con $P=\V$ el miembro izquierdo exige $Q$ y $\neg Q$, luego vale $\F=\neg P$; con $P=\F$ ambos condicionales son verdaderos y $\neg P=\V$.
\end{solucion}

\begin{solucion}{eje:1-validez}
\datos cuatro razonamientos.
\planteamiento identificar la forma lógica y compararla con la \cref{prop:inferencia} y la \cref{obs:1-errores}.
\pasos
\begin{enumerate}[label=(\alph*)]
  \item $P\Rightarrow Q$, $Q$; luego $P$. Es la falacia de \textbf{afirmar el consecuente}: el error de validación puede subir por otras causas. \textbf{No válido.}
  \item $P\Rightarrow Q$ con $P$: «$n$ primo mayor que~2», $Q$: «$n$ impar»; premisa $\neg Q$; conclusión $\neg P$. \textbf{Modus tollens, válido.}
  \item $P\vee Q$, $\neg P$; luego $Q$. \textbf{Silogismo disyuntivo, válido.}
  \item $P\Rightarrow Q$, $\neg P$; luego $\neg Q$. Falacia de \textbf{negar el antecedente}; contraejemplo: $x=-3\le 2$ y $x^{2}=9>4$. \textbf{No válido.}
\end{enumerate}
\resultado válidos (b) y (c); no válidos (a) y (d).
\comprobacion en (a) y (d), la asignación $P=\F$, $Q=\V$ hace verdaderas las premisas y falsa la conclusión.
\end{solucion}

\begin{solucion}{eje:1-contar}
\datos $n$ variables.
\planteamiento una tabla asigna un valor a cada una de las $2^{n}$ filas (\cref{prop:filas}).
\pasos
\begin{enumerate}[label=\arabic*.]
  \item Cada una de las $2^{n}$ filas recibe $\V$ o $\F$ de forma independiente: $2^{2^{n}}$ tablas posibles.
  \item $n=1$: $2^{2}=4$; $n=2$: $2^{4}=16$; $n=3$: $2^{8}=256$.
  \item Para $n=2$, solo una tabla es toda $\V$ (tautología) y solo una es toda $\F$ (contradicción).
  \item Las cuatro funciones de una variable: constante $\V$ ($P\vee\neg P$), constante $\F$ ($P\wedge\neg P$), identidad ($P$) y negación ($\neg P$).
\end{enumerate}
\resultado (a) $2^{2^{n}}$: \textbf{4, 16, 256}; (b) \textbf{una tautología y una contradicción}; (c) $P\vee\neg P$, $P\wedge\neg P$, $P$, $\neg P$.
\comprobacion la \cref{tab:conectivas} y el \cref{eje:1-xor} muestran ya 6 de las 16 funciones de dos variables: $\wedge$, $\vee$, $\Rightarrow$, $\Leftrightarrow$, $\xor$ y la proyección $P$.
\end{solucion}

\begin{solucion}{eje:1-leyes}
\datos leyes (iii), (vi) y (vii).
\planteamiento tablas de verdad; para (iii) basta razonar por casos según $P$.
\pasos
\begin{enumerate}[label=\arabic*.]
  \item (iii), primera ley. Si $P=\F$, ambos miembros valen $\F$. Si $P=\V$, el izquierdo vale $Q\vee R$ y el derecho $(\V\wedge Q)\vee(\V\wedge R)=Q\vee R$. Coinciden en las 8 filas.
  \item (iii), segunda ley. Si $P=\V$, ambos valen $\V$. Si $P=\F$, el izquierdo vale $Q\wedge R$ y el derecho $(\F\vee Q)\wedge(\F\vee R)=Q\wedge R$.
  \item (vi). $\neg(P\Rightarrow Q)$ es verdadera solo en la fila $P=\V,Q=\F$ (\cref{tab:conectivas}); $P\wedge\neg Q$ también. En las demás filas ambas son falsas.
  \item (vii). $(P\Rightarrow Q)\wedge(Q\Rightarrow P)$: en VV vale $\V\wedge\V=\V$; en VF, $\F\wedge\V=\F$; en FV, $\V\wedge\F=\F$; en FF, $\V\wedge\V=\V$. Es la columna de $P\Leftrightarrow Q$.
\end{enumerate}
\resultado las tres leyes \textbf{se cumplen} en todas las asignaciones.
\comprobacion (vi) se deduce también de (iv) y De Morgan: $\neg(\neg P\vee Q)\eqlog\neg\neg P\wedge\neg Q\eqlog P\wedge\neg Q$.
\end{solucion}

\begin{solucion}{eje:1-demorgan-n}
\datos $n\ge 2$ proposiciones $P_1,\dots,P_n$.
\planteamiento dos fórmulas son equivalentes si son verdaderas exactamente en las mismas asignaciones; se describe ese conjunto para cada miembro.
\pasos
\begin{enumerate}[label=\arabic*.]
  \item Sea una asignación cualquiera. Por definición de la conjunción (aplicada repetidamente), $P_1\wedge\dots\wedge P_n$ es verdadera si y solo si todas las $P_i$ son verdaderas.
  \item Por tanto, $\neg(P_1\wedge\dots\wedge P_n)$ es verdadera si y solo si \emph{no} todas las $P_i$ son verdaderas, es decir, si y solo si existe un índice $i$ con $P_i$ falsa.
  \item Por otra parte, $\neg P_1\vee\dots\vee\neg P_n$ es verdadera si y solo si alguna $\neg P_i$ es verdadera, es decir, si y solo si existe $i$ con $P_i$ falsa.
  \item Las dos condiciones de los pasos 2 y 3 son la misma, luego ambas fórmulas tienen el mismo valor en cada asignación: son equivalentes.
  \item Para la segunda ley, se aplica la primera a las proposiciones $\neg P_1,\dots,\neg P_n$: $\neg(\neg P_1\wedge\dots\wedge\neg P_n)\eqlog\neg\neg P_1\vee\dots\vee\neg\neg P_n\eqlog P_1\vee\dots\vee P_n$. Negando ambos miembros y usando la doble negación se obtiene $\neg P_1\wedge\dots\wedge\neg P_n\eqlog\neg(P_1\vee\dots\vee P_n)$. $\blacksquare$
\end{enumerate}
\resultado \textbf{las leyes de De Morgan valen para cualquier número finito de proposiciones}.
\comprobacion con $n=2$ se recupera la \cref{prop:leyes}(ii). Una prueba alternativa es por inducción sobre $n$ (\cref{tema:5}), escribiendo $P_1\wedge\dots\wedge P_{n+1}=(P_1\wedge\dots\wedge P_n)\wedge P_{n+1}$.
\end{solucion}

\begin{solucion}{eje:1-completitud}
\datos una tabla de verdad $T$ con variables $P_1,\dots,P_n$.
\planteamiento construir la forma normal disyuntiva y eliminar $\vee$ con De Morgan.
\pasos
\begin{enumerate}[label=\arabic*.]
  \item Para cada asignación $a=(a_1,\dots,a_n)$ se define el \emph{literal} $L_i^{a}=P_i$ si $a_i=\V$ y $L_i^{a}=\neg P_i$ si $a_i=\F$, y la conjunción $C_a=L_1^{a}\wedge\dots\wedge L_n^{a}$. Por construcción, $C_a$ es verdadera en la asignación $a$ y falsa en todas las demás, pues en cualquier otra asignación algún literal es falso.
  \item Sea $A$ el conjunto de asignaciones en las que $T$ vale $\V$. Si $A=\varnothing$, la fórmula $P_1\wedge\neg P_1$ tiene la tabla $T$. Si $A\neq\varnothing$, sea $\Phi=\bigvee_{a\in A}C_a$.
  \item $\Phi$ tiene la tabla $T$: en una asignación $b\in A$, el término $C_b$ es verdadero y $\Phi$ también; en una asignación $b\notin A$, todos los $C_a$ con $a\in A$ son falsos, porque $a\neq b$, y $\Phi$ es falsa.
  \item Se elimina $\vee$: por De Morgan generalizado (\cref{eje:1-demorgan-n}) y la doble negación, $\bigvee_{a\in A}C_a\eqlog\neg\bigwedge_{a\in A}\neg C_a$, que solo contiene $\neg$ y $\wedge$.
  \item Por tanto, toda tabla de verdad es la de una fórmula con solo $\neg$ y $\wedge$. $\blacksquare$
\end{enumerate}
\resultado \textbf{$\{\neg,\wedge\}$ es un conjunto completo de conectivas}.
\comprobacion para $T=P\xor Q$, $A=\{(\V,\F),(\F,\V)\}$ y $\Phi=(P\wedge\neg Q)\vee(\neg P\wedge Q)$, que coincide con la tabla del \cref{eje:1-xor}; su forma con $\neg,\wedge$ es $\neg\bigl(\neg(P\wedge\neg Q)\wedge\neg(\neg P\wedge Q)\bigr)$.
\end{solucion}

\begin{solucion}{eje:1-bicond-asoc}
\datos $L=(P\Leftrightarrow Q)\Leftrightarrow R$, $D=P\Leftrightarrow(Q\Leftrightarrow R)$.
\planteamiento caracterizar cuándo es verdadero un bicondicional en términos del número de falsas.
\pasos
\begin{enumerate}[label=\arabic*.]
  \item \emph{Lema.} $A\Leftrightarrow B$ es verdadero si y solo si el número de falsas entre $A$ y $B$ es par (0 o 2). En efecto, es verdadero cuando ambas tienen el mismo valor, es decir, con 0 falsas (ambas $\V$) o 2 falsas (ambas $\F$).
  \item Sea $f$ el número de falsas entre $P,Q,R$. Si $R=\V$, $L$ tiene el valor de $P\Leftrightarrow Q$, verdadero si y solo si el número de falsas entre $P,Q$ es par, es decir (como $R$ no aporta falsas), si y solo si $f$ es par.
  \item Si $R=\F$, $L$ es verdadera si y solo si $P\Leftrightarrow Q$ es falsa, es decir, si el número de falsas entre $P,Q$ es impar; sumando la falsa $R$, si y solo si $f$ es par.
  \item El mismo argumento, separando según el valor de $P$, prueba que $D$ es verdadera si y solo si $f$ es par.
  \item Ambas fórmulas son verdaderas exactamente en las asignaciones con $f$ par; luego $L\eqlog D$. $\blacksquare$
  \item El condicional no es asociativo: con $P=Q=R=\F$, $(P\Rightarrow Q)\Rightarrow R=\V\Rightarrow\F=\F$, mientras que $P\Rightarrow(Q\Rightarrow R)=\F\Rightarrow\V=\V$.
\end{enumerate}
\resultado \textbf{$\Leftrightarrow$ es asociativo; $\Rightarrow$ no lo es} (asignación $P=Q=R=\F$).
\comprobacion fila $P=\V$, $Q=\F$, $R=\F$ ($f=2$, par): $L=(\F)\Leftrightarrow\F=\V$ y $D=\V\Leftrightarrow(\F\Leftrightarrow\F)=\V\Leftrightarrow\V=\V$. Se han comprobado además las 8 filas por ordenador.
\end{solucion}

% ---------------------------------------------------------------------
\section{Tema 2}

\begin{solucion}{eje:2-valores}
\datos seis enunciados cuantificados; $\N=\{1,2,\dots\}$.
\planteamiento un universal se prueba en general y se refuta con un contraejemplo; un existencial se prueba con un testigo.
\pasos
\begin{enumerate}[label=(\alph*)]
  \item \textbf{V}: el cuadrado de un real es no negativo.
  \item \textbf{F}: $\sqrt2\notin\Q$ (\cref{teo:raiz2}), y las únicas soluciones reales son $\pm\sqrt2$.
  \item \textbf{V}: testigo $x=\sqrt2$.
  \item \textbf{F}: contraejemplo $n=1$, pues no hay $m\in\N$ con $m<1$.
  \item \textbf{V}: testigo $m=1$, que cumple $1\le n$ para todo $n\in\N$.
  \item \textbf{V}: universal sobre el conjunto vacío (vacuamente verdadero).
\end{enumerate}
\resultado V, F, V, F, V, V.
\comprobacion (d) y (e) muestran que en $\N$ hay mínimo; con $\Z$ en lugar de $\N$ los valores se invertirían.
\end{solucion}

\begin{solucion}{eje:2-simbolos}
\datos cuatro enunciados.
\planteamiento identificar universo, cuantificadores y forma (universal condicional, existencial conjuntiva).
\pasos
\begin{enumerate}[label=(\alph*)]
  \item $\forall n\in\Z\ \bigl((n\text{ par}\wedge n>2)\Rightarrow\exists p\,\exists q\,(p,q\text{ primos}\wedge n=p+q)\bigr)$.
  \item $\exists x\in\R\ \forall n\in\Z\ (x>n)$. \textbf{Falso}: dado $x$, el entero $n=\ceil{x}$ cumple $n\ge x$, luego no es cierto que $x>n$.
  \item $\forall n\in\Z\ \bigl(n^{2}\not\equiv 2\pmod{10}\bigr)$, o $\neg\exists n\in\Z\ (n^{2}\equiv 2\pmod{10})$.
  \item $\neg\exists x\in\R\ (x^{2}+1=0)$, equivalente a $\forall x\in\R\ (x^{2}+1\neq 0)$.
\end{enumerate}
\resultado las fórmulas anteriores; (b) es falso.
\comprobacion (c) es verdadero: los cuadrados de $0,\dots,9$ terminan en $0,1,4,9,6,5,6,9,4,1$, y la última cifra de $n^{2}$ solo depende de la de $n$ (\cref{prop:congruencias}).
\end{solucion}

\begin{solucion}{eje:2-lenguaje}
\datos cuatro fórmulas.
\planteamiento leer de izquierda a derecha, teniendo presente que una variable cuantificada después puede depender de las anteriores.
\pasos
\begin{enumerate}[label=(\alph*)]
  \item «Para todo real hay otro mayor». \textbf{V}: $y=x+1$.
  \item «Hay un real mayor que todos los reales». \textbf{F}: tomando $x=y$ sería $y>y$.
  \item «Todo real no nulo tiene inverso». \textbf{V}: $y=1/x$.
  \item «Hay un entero que, multiplicado por cualquier entero, lo deja igual». \textbf{V}: $x=1$.
\end{enumerate}
\resultado V, F, V, V.
\comprobacion (a) y (b) son el par $\forall\exists$/$\exists\forall$ de la \cref{prop:orden} con $P(x,y)$: «$y>x$».
\end{solucion}

\begin{solucion}{eje:2-negar}
\datos cuatro enunciados.
\planteamiento aplicar el \cref{teo:negacion}, De Morgan y $\neg(A\Rightarrow B)\eqlog A\wedge\neg B$.
\pasos
\begin{enumerate}[label=(\alph*)]
  \item Negación: $\exists x\in\R\ \forall y\in\R\ (y\le x\vee y\ge x+1)$. El original es \textbf{verdadero}: $y=x+\tfrac12$.
  \item Negación: $\forall n\in\N\ \exists m\in\N\ (n\nodivide m)$. El original es \textbf{verdadero}: $n=1$ divide a todo $m$.
  \item Negación: $\exists x\in\R\ (x^{2}>4\wedge x\le 2)$. La \textbf{negación es verdadera}: $x=-3$.
  \item Negación: $\exists n\in\Z\ (n\text{ primo}\wedge n\text{ par}\wedge n\neq 2)$. El original es \textbf{verdadero}: un primo par es divisible por~2, luego igual a~2.
\end{enumerate}
\resultado verdaderos: original en (a), (b), (d); negación en (c).
\comprobacion en (d), $\neg(A\vee B)\eqlog\neg A\wedge\neg B$ convierte «impar o igual a~2» en «par y distinto de~2».
\end{solucion}

\begin{solucion}{eje:2-continuidad}
\datos definición $\varepsilon$–$\delta$; $f(x)=0$ si $x<0$, $f(x)=1$ si $x\ge 0$; $a=0$.
\planteamiento negar como en el \cref{ej:negar-limite} y exhibir $\varepsilon$ y, para cada $\delta$, un $x$.
\pasos
\begin{enumerate}[label=\arabic*.]
  \item Negación: $\exists\varepsilon>0\ \forall\delta>0\ \exists x\in\R\ \bigl(\abs{x-a}<\delta\wedge\abs{f(x)-f(a)}\ge\varepsilon\bigr)$.
  \item Se toma $\varepsilon=\tfrac12$. Sea $\delta>0$ arbitrario y sea $x=-\delta/2$.
  \item Entonces $\abs{x-0}=\delta/2<\delta$ y, como $x<0$, $\abs{f(x)-f(0)}=\abs{0-1}=1\ge\tfrac12$.
  \item Como $\delta$ era arbitrario, se cumple la negación: $f$ no es continua en $0$. $\blacksquare$
\end{enumerate}
\resultado negación del paso~1; \textbf{$f$ no es continua en $0$} (con $\varepsilon=\tfrac12$ y $x=-\delta/2$).
\comprobacion cualquier $\varepsilon\in(0,1]$ sirve; con $\varepsilon>1$ no, porque $\abs{f(x)-f(0)}\le 1$.
\end{solucion}

\begin{solucion}{eje:2-orden}
\datos cuatro predicados y universos.
\planteamiento para $\forall x\exists y$, dar $y$ en función de $x$; para $\exists y\forall x$, un $y$ fijo o refutarlo.
\pasos
\begin{enumerate}[label=(\alph*)]
  \item $\N_0$, $x+y=0$. $\forall\exists$: \textbf{F} ($x=1$ exigiría $y=-1\notin\N_0$). $\exists\forall$: \textbf{F} (implica el anterior, \cref{prop:orden}).
  \item $\Q\setminus\{0\}$, $xy=1$. $\forall\exists$: \textbf{V} ($y=1/x$). $\exists\forall$: \textbf{F} ($x=1$ da $y=1$ y $x=2$ da $y=\tfrac12$).
  \item $\N$, $x\le y$. $\forall\exists$: \textbf{V} ($y=x$). $\exists\forall$: \textbf{F} (sería un máximo de $\N$; $x=y+1$ lo refuta).
  \item $\N$, $y\le x$. $\forall\exists$: \textbf{V} ($y=1$). $\exists\forall$: \textbf{V} ($y=1$, el mínimo de $\N$).
\end{enumerate}
\resultado (a) F, F; (b) V, F; (c) V, F; (d) V, V.
\comprobacion en ningún caso se tiene $\exists\forall$ verdadero y $\forall\exists$ falso, como exige la \cref{prop:orden}.
\end{solucion}

\begin{solucion}{eje:2-distribuir}
\datos cuatro posibles equivalencias.
\planteamiento para las falsas, universo $\Z$, $P(x)$: «$x$ par», $Q(x)$: «$x$ impar».
\pasos
\begin{enumerate}[label=(\alph*)]
  \item \textbf{Cierta} (demostración en el \cref{eje:2-distribuir-dem}).
  \item \textbf{Falsa}: $\forall x\,(P(x)\vee Q(x))$ es verdadera (todo entero es par o impar), pero $\forall x\,P(x)$ y $\forall x\,Q(x)$ son falsas.
  \item \textbf{Cierta} (\cref{eje:2-distribuir-dem}).
  \item \textbf{Falsa}: $\exists x\,P(x)$ y $\exists x\,Q(x)$ son verdaderas (testigos 2 y 3), pero $\exists x\,(P(x)\wedge Q(x))$ es falsa.
\end{enumerate}
\resultado $\forall$ distribuye sobre $\wedge$ y $\exists$ sobre $\vee$; no al revés.
\comprobacion (b) y (d) son negación una de la otra: aplicando el \cref{teo:negacion} a (b) con $\neg P$, $\neg Q$ se obtiene (d), luego una falla si y solo si falla la otra.
\end{solucion}

\begin{solucion}{eje:2-distribuir-dem}
\datos predicados $P,Q$ sobre un universo $U$.
\planteamiento probar cada bicondicional por sus dos implicaciones, de forma directa.
\pasos
\begin{enumerate}[label=\arabic*.]
  \item (a, $\Rightarrow$) Supóngase $\forall x\,(P(x)\wedge Q(x))$. Sea $a\in U$ arbitrario. Entonces $P(a)\wedge Q(a)$, en particular $P(a)$; como $a$ es arbitrario, $\forall x\,P(x)$. Del mismo modo, $\forall x\,Q(x)$. Luego $\forall x\,P(x)\wedge\forall x\,Q(x)$.
  \item (a, $\Leftarrow$) Supóngase $\forall x\,P(x)$ y $\forall x\,Q(x)$. Sea $a\in U$ arbitrario. Por la primera, $P(a)$; por la segunda, $Q(a)$; luego $P(a)\wedge Q(a)$. Como $a$ es arbitrario, $\forall x\,(P(x)\wedge Q(x))$.
  \item (c, $\Rightarrow$) Supóngase $\exists x\,(P(x)\vee Q(x))$ y sea $a$ un testigo: $P(a)\vee Q(a)$. Si $P(a)$, entonces $\exists x\,P(x)$; si $Q(a)$, entonces $\exists x\,Q(x)$. En ambos casos, $\exists x\,P(x)\vee\exists x\,Q(x)$.
  \item (c, $\Leftarrow$) Supóngase $\exists x\,P(x)\vee\exists x\,Q(x)$. Si $\exists x\,P(x)$ con testigo $a$, entonces $P(a)\vee Q(a)$ y $a$ es testigo de $\exists x\,(P(x)\vee Q(x))$. El otro caso es análogo.
  \item (Implicación adicional) Supóngase $\forall x\,P(x)\vee\forall x\,Q(x)$ y sea $a\in U$. Si vale $\forall x\,P(x)$, entonces $P(a)$, luego $P(a)\vee Q(a)$; en el otro caso, $Q(a)$ y también $P(a)\vee Q(a)$. Como $a$ es arbitrario, $\forall x\,(P(x)\vee Q(x))$. $\blacksquare$
\end{enumerate}
\resultado \textbf{(a) y (c) son equivalencias; $\forall x\,P\vee\forall x\,Q\Rightarrow\forall x\,(P\vee Q)$ siempre}, aunque su recíproca falla (\cref{eje:2-distribuir}(b)).
\comprobacion los pasos 3, 4 y 5 son demostraciones por casos (\cref{def:casos}); la palabra «arbitrario» marca dónde se demuestra un universal.
\end{solucion}

\begin{solucion}{eje:2-vacio}
\datos $U\neq\varnothing$; predicado $P$.
\planteamiento demostración directa, usando un elemento de $U$ como testigo.
\pasos
\begin{enumerate}[label=\arabic*.]
  \item Supóngase $\forall x\in U\ P(x)$. Como $U$ es no vacío, existe algún $a\in U$.
  \item Por la hipótesis aplicada a $a$, $P(a)$ es verdadera.
  \item Luego $a$ es testigo de $\exists x\in U\ P(x)$. $\blacksquare$
  \item Si $U=\varnothing$, $\forall x\,P(x)$ es vacuamente verdadero y $\exists x\,P(x)$ es falso; el condicional sería $\V\Rightarrow\F$, falso. El paso~1 es el que falla: no hay ningún $a$.
\end{enumerate}
\resultado \textbf{la implicación vale si $U\neq\varnothing$} y falla para $U=\varnothing$.
\comprobacion la \cref{def:predicado} exige universos no vacíos precisamente para evitar este caso.
\end{solucion}

\begin{solucion}{eje:2-unicidad}
\datos \cref{def:unicidad}.
\planteamiento separar existencia y unicidad; negar con el \cref{teo:negacion}.
\pasos
\begin{enumerate}[label=(\alph*)]
  \item $\exists x\,P(x)\ \wedge\ \forall x\,\forall y\,\bigl((P(x)\wedge P(y))\Rightarrow x=y\bigr)$.
  \item Negación: $\neg\exists x\,P(x)\vee\exists x\,\exists y\,(P(x)\wedge P(y)\wedge x\neq y)$, es decir, $\forall x\,\neg P(x)\ \vee\ \exists x\,\exists y\,(P(x)\wedge P(y)\wedge x\neq y)$. Puede fallar porque no haya ningún elemento con la propiedad o porque haya al menos dos distintos.
  \item En $\R$, $x^{2}=4$ tiene dos soluciones, $2$ y $-2$: \textbf{falso}. En $\N$, solo $x=2$: \textbf{verdadero}.
\end{enumerate}
\resultado las fórmulas de (a) y (b); (c) falso en $\R$ y verdadero en $\N$.
\comprobacion la fórmula de (a) es equivalente a la de la \cref{def:unicidad}: si $a$ cumple $P(a)\wedge\forall y\,(P(y)\Rightarrow y=a)$, dos elementos $x,y$ con la propiedad son ambos iguales a $a$.
\end{solucion}

\begin{solucion}{eje:2-acotada}
\datos definiciones de acotada y creciente; $f(x)=x^{2}$.
\planteamiento negar y dar testigos.
\pasos
\begin{enumerate}[label=\arabic*.]
  \item No acotada: $\forall M\in\R\ \exists x\in\R\ (\abs{f(x)}>M)$.
  \item No creciente: $\exists x\,\exists y\ (x<y\wedge f(x)>f(y))$.
  \item $f$ no es acotada: sea $M\in\R$ arbitrario y $x=\abs{M}+1$. Entonces $x\ge 1$, luego $x^{2}\ge x=\abs{M}+1>M$.
  \item $f$ no es creciente: $x=-2<y=0$ y $f(-2)=4>0=f(0)$. $\blacksquare$
\end{enumerate}
\resultado negaciones de los pasos 1 y 2; $x^{2}$ \textbf{no es acotada ni creciente}.
\comprobacion la negación de «no acotada» requiere un $x$ \emph{para cada} $M$: el $x$ del paso~3 depende de $M$, como debe ser.
\end{solucion}

\begin{solucion}{eje:2-arquimedes}
\datos propiedad arquimediana.
\planteamiento (a) dado $\varepsilon$, elegir $N>1/\varepsilon$; (b) negar y exhibir, para cada $N$, un $\varepsilon$ y un $n$.
\pasos
\begin{enumerate}[label=\arabic*.]
  \item (a) Sea $\varepsilon>0$. Por la propiedad arquimediana aplicada a $x=1/\varepsilon$, existe $N\in\N$ con $N>1/\varepsilon$. Sea $n\in\N$ con $n\ge N$. Entonces $n>1/\varepsilon$, y como $n>0$ y $\varepsilon>0$, $1/n<\varepsilon$. Esto prueba (a).
  \item (b) La negación es $\forall N\in\N\ \exists\varepsilon>0\ \exists n\in\N\ (n\ge N\wedge 1/n\ge\varepsilon)$. Sea $N\in\N$ arbitrario. Se toman $n=N$ y $\varepsilon=1/N>0$. Entonces $n\ge N$ y $1/n=\varepsilon\ge\varepsilon$.
  \item Se ha probado la negación de (b), luego (b) es falso. $\blacksquare$
\end{enumerate}
\resultado \textbf{(a) verdadero; (b) falso}. En (a), $N$ depende de $\varepsilon$; (b) pediría un $N$ válido para todas las tolerancias.
\comprobacion con $\varepsilon=0{,}01$, (a) da $N=101$, y $1/101\approx 0{,}0099<0{,}01$.
\end{solucion}

\begin{solucion}{eje:2-critica}
\datos «todo natural es impar» y la prueba propuesta.
\planteamiento aplicar el \cref{teo:negacion}.
\pasos
\begin{enumerate}[label=\arabic*.]
  \item La negación de $\forall n\in\N\ (n\text{ impar})$ es $\exists n\in\N\ (n\text{ no impar})$, es decir, «algún natural es par», no «todos los naturales son pares».
  \item La prueba refuta un enunciado («todos son pares») que no es la negación del original; la contradicción obtenida no dice nada sobre él.
  \item El enunciado es \textbf{falso}: $2\in\N$ es par.
\end{enumerate}
\resultado error de \textbf{negación} (\cref{obs:2-errores}(a)); el teorema es falso, contraejemplo $n=2$.
\comprobacion el mismo «argumento» probaría, cambiando 137 por 2, que todo natural es par: señal de que no prueba nada.
\end{solucion}

\begin{solucion}{eje:2-finito}
\datos $U=\{a_1,\dots,a_n\}$.
\planteamiento traducir la \cref{def:cuantificadores} y usar el \cref{eje:1-demorgan-n}.
\pasos
\begin{enumerate}[label=\arabic*.]
  \item $\forall x\,P(x)$ es verdadera si y solo si $P(a)$ lo es para cada $a\in U$, es decir, si y solo si $P(a_1),\dots,P(a_n)$ son todas verdaderas, que es la condición de verdad de $P(a_1)\wedge\dots\wedge P(a_n)$.
  \item $\exists x\,P(x)$ es verdadera si y solo si alguna $P(a_i)$ lo es, condición de verdad de $P(a_1)\vee\dots\vee P(a_n)$.
  \item Entonces $\neg\forall x\,P(x)\eqlog\neg\bigl(P(a_1)\wedge\dots\wedge P(a_n)\bigr)\eqlog\neg P(a_1)\vee\dots\vee\neg P(a_n)\eqlog\exists x\,\neg P(x)$, usando los pasos 1 y 2 (este último con el predicado $\neg P$) y De Morgan generalizado.
  \item La segunda ley se obtiene igual con la otra ley de De Morgan. $\blacksquare$
\end{enumerate}
\resultado en universos finitos, $\forall$ es una conjunción, $\exists$ una disyunción, y el \cref{teo:negacion} es De Morgan.
\comprobacion la demostración del \cref{teo:negacion} en el texto no usa la finitud, y vale para universos infinitos como $\N$ o $\R$.
\end{solucion}

% ---------------------------------------------------------------------
\section{Tema 3}

\begin{solucion}{eje:3-hipotesis}
\datos cuatro condicionales.
\planteamiento separar hipótesis y tesis, aplicar la \cref{def:reciproco} y valorar qué punto de partida es más manejable.
\pasos
\begin{enumerate}[label=(\alph*)]
  \item Hipótesis: $a\divides b$ y $b\divides c$; tesis: $a\divides c$. Contrarrecíproco: si $a\nodivide c$, entonces $a\nodivide b$ o $b\nodivide c$. \textbf{Directa}: las hipótesis dan igualdades $b=ad$, $c=be$ (\cref{prop:divisibilidad}).
  \item Hipótesis: $n^{2}$ impar; tesis: $n$ impar. Contrarrecíproco: si $n$ es par, $n^{2}$ es par. \textbf{Contrarrecíproco}: partir de $n=2k$ es inmediato; partir de $n^{2}=2k+1$ no.
  \item Hipótesis: $x^{2}+x<0$; tesis: $x<0$. Contrarrecíproco: si $x\ge 0$, entonces $x^{2}+x\ge 0$. \textbf{Contrarrecíproco}: suma de no negativos.
  \item Hipótesis: $xy$ irracional; tesis: $x$ o $y$ irracional. Contrarrecíproco: si $x$ e $y$ son racionales, $xy$ es racional. \textbf{Contrarrecíproco}: «irracional» es una negación y su contrario, «racional», da una fracción con la que operar.
\end{enumerate}
\resultado directa en (a); contrarrecíproco en (b), (c) y (d).
\comprobacion en (d) la negación de «$x$ o $y$ irracional» es «$x$ e $y$ racionales» por De Morgan.
\end{solucion}

\begin{solucion}{eje:3-esquema}
\datos la prueba propuesta.
\planteamiento comparar con la \cref{def:directa}.
\pasos
\begin{enumerate}[label=\arabic*.]
  \item Defecto 1: parte de la tesis ($n^{2}=4k$) en lugar de la hipótesis ($n$ par): demuestra, si acaso, el recíproco.
  \item Defecto 2: $2\sqrt{k}$ no es necesariamente un entero, de modo que «$n=2\sqrt k$ es par» no está justificado.
  \item Redacción correcta: «Sea $n$ un entero par. Por definición existe $k\in\Z$ con $n=2k$. Entonces $n^{2}=4k^{2}$, y como $k^{2}\in\Z$, $n^{2}$ es múltiplo de~4. $\blacksquare$»
\end{enumerate}
\resultado error de dirección (se supone la tesis) y paso no justificado; prueba correcta en el paso~3.
\comprobacion con $n=6$: $n^{2}=36=4\cdot 9$.
\end{solucion}

\begin{solucion}{eje:3-cubo}
\datos $x\in\Z$ impar.
\planteamiento demostración directa: desplegar la definición y llegar a la forma $2m+1$.
\pasos
\begin{enumerate}[label=\arabic*.]
  \item Sea $x$ un entero impar. Por la \cref{def:paridad}, existe $k\in\Z$ con $x=2k+1$.
  \item Entonces $x^{3}=(2k+1)^{3}=8k^{3}+12k^{2}+6k+1=2(4k^{3}+6k^{2}+3k)+1$.
  \item Sea $m=4k^{3}+6k^{2}+3k$, que es entero por ser suma de productos de enteros. Entonces $x^{3}=2m+1$, luego $x^{3}$ es impar. $\blacksquare$
\end{enumerate}
\resultado \textbf{si $x$ es impar, $x^{3}$ es impar.}
\comprobacion $x=3$: $k=1$, $m=4+6+3=13$, y $2\cdot 13+1=27=3^{3}$.
\end{solucion}

\begin{solucion}{eje:3-polinomio}
\datos $a\in\Z$ impar.
\planteamiento directa, como en el \cref{eje:3-cubo}.
\pasos
\begin{enumerate}[label=\arabic*.]
  \item Sea $a$ impar; existe $k\in\Z$ con $a=2k+1$.
  \item $a^{2}+3a+5=(4k^{2}+4k+1)+(6k+3)+5=4k^{2}+10k+9=2(2k^{2}+5k+4)+1$.
  \item Con $m=2k^{2}+5k+4\in\Z$, $a^{2}+3a+5=2m+1$ es impar. $\blacksquare$
\end{enumerate}
\resultado \textbf{$a^{2}+3a+5$ es impar si $a$ es impar.}
\comprobacion $a=1$: $1+3+5=9$ impar, y la fórmula da $k=0$, $m=4$, $2\cdot 4+1=9$. Otra vía: $a^{2}$ y $3a$ son impares, su suma es par, y par más~5 es impar.
\end{solucion}

\begin{solucion}{eje:3-cuadrados}
\datos $a,b\in\Z$ con $a\divides b$.
\planteamiento directa con la \cref{def:divide}.
\pasos
\begin{enumerate}[label=\arabic*.]
  \item Supóngase $a\divides b$. Existe $c\in\Z$ con $b=ac$.
  \item Elevando al cuadrado, $b^{2}=a^{2}c^{2}$.
  \item Como $c^{2}\in\Z$, la igualdad muestra que $a^{2}\divides b^{2}$. $\blacksquare$
\end{enumerate}
\resultado \textbf{$a\divides b\Rightarrow a^{2}\divides b^{2}$.}
\comprobacion $3\divides 12$ y $9\divides 144=9\cdot 16$. El recíproco también es cierto, pero su demostración requiere la factorización única en primos.
\end{solucion}

\begin{solucion}{eje:3-producto}
\datos $a\divides b$, $c\divides d$.
\planteamiento directa; usar letras distintas para los dos cocientes.
\pasos
\begin{enumerate}[label=\arabic*.]
  \item Supóngase $a\divides b$ y $c\divides d$. Existen $k,\ell\in\Z$ con $b=ak$ y $d=c\ell$.
  \item Multiplicando, $bd=(ak)(c\ell)=(ac)(k\ell)$.
  \item Como $k\ell\in\Z$, se concluye $ac\divides bd$. $\blacksquare$
\end{enumerate}
\resultado \textbf{$ac\divides bd$.}
\comprobacion $2\divides 6$ y $5\divides 15$ dan $10\divides 90$; en efecto, $90=10\cdot 9$ y $k\ell=3\cdot 3=9$.
\end{solucion}

\begin{solucion}{eje:3-desigualdad}
\datos $x\in\R$, $0<x<4$.
\planteamiento buscar hacia atrás: $\frac{4}{x(4-x)}\ge 1\iff 4\ge x(4-x)$ (por ser $x(4-x)>0$) $\iff x^{2}-4x+4\ge 0\iff(x-2)^{2}\ge 0$. Después se redacta hacia delante.
\pasos
\begin{enumerate}[label=\arabic*.]
  \item Sea $x\in\R$ con $0<x<4$. Como el cuadrado de un real es no negativo, $0\le(x-2)^{2}=x^{2}-4x+4$.
  \item Sumando $4x-x^{2}$ a ambos miembros, $x(4-x)=4x-x^{2}\le 4$.
  \item De $0<x$ y $x<4$ se obtiene $x>0$ y $4-x>0$, luego $x(4-x)>0$.
  \item Dividiendo la desigualdad del paso~2 por el número positivo $x(4-x)$, se obtiene $1\le\dfrac{4}{x(4-x)}$. $\blacksquare$
\end{enumerate}
\resultado \textbf{$\dfrac{4}{x(4-x)}\ge 1$}, con igualdad solo si $x=2$.
\comprobacion $x=1$: $4/3\ge 1$; $x=2$: $4/4=1$. La hipótesis $0<x<4$ se usa en el paso~3: sin ella el denominador podría ser negativo o nulo.
\end{solucion}

\begin{solucion}{eje:3-diferencia}
\datos $n\in\Z$ impar.
\planteamiento buscar $a,b$ con $a^{2}-b^{2}=(a-b)(a+b)=n$; con $a-b=1$ basta $a+b=n$.
\pasos
\begin{enumerate}[label=\arabic*.]
  \item Sea $n$ un entero impar; existe $k\in\Z$ con $n=2k+1$.
  \item Sean $a=k+1$ y $b=k$, que son enteros.
  \item Entonces $a^{2}-b^{2}=(k+1)^{2}-k^{2}=k^{2}+2k+1-k^{2}=2k+1=n$.
  \item Por tanto, $n$ es diferencia de los cuadrados de dos enteros. $\blacksquare$
\end{enumerate}
\resultado \textbf{$2k+1=(k+1)^{2}-k^{2}$.}
\comprobacion $n=9$: $k=4$, $25-16=9$; $n=-3$: $k=-2$, $(-1)^{2}-(-2)^{2}=1-4=-3$.
\end{solucion}

\begin{solucion}{eje:3-contra-cubo}
\datos $n\in\Z$.
\planteamiento contrarrecíproco: si $n$ es impar, $n^{3}+5$ es par.
\pasos
\begin{enumerate}[label=\arabic*.]
  \item Se demuestra el contrarrecíproco. Sea $n$ un entero no par, es decir, impar.
  \item Por el \cref{eje:3-cubo}, $n^{3}$ es impar: $n^{3}=2m+1$ con $m\in\Z$.
  \item Entonces $n^{3}+5=2m+6=2(m+3)$, que es par, es decir, no es impar.
  \item Se ha probado «si $n$ no es par, $n^{3}+5$ no es impar», que equivale al enunciado. $\blacksquare$
\end{enumerate}
\resultado \textbf{si $n^{3}+5$ es impar, $n$ es par.}
\comprobacion $n=2$: $13$ impar y $n$ par; $n=1$: $6$ par, coherente con el contrarrecíproco.
\end{solucion}

\begin{solucion}{eje:3-pitagoras}
\datos $a,b,c\in\Z$ con $a^{2}+b^{2}=c^{2}$.
\planteamiento contrarrecíproco: si $a,b,c$ son todos impares, $a^{2}+b^{2}\neq c^{2}$.
\pasos
\begin{enumerate}[label=\arabic*.]
  \item La negación de «al menos uno es par» es «los tres son impares». Se demuestra: si $a$, $b$ y $c$ son impares, entonces $a^{2}+b^{2}\neq c^{2}$.
  \item Sean $a,b,c$ impares. Por la \cref{prop:impar-cuadrado}, $a^{2}$, $b^{2}$ y $c^{2}$ son impares: existen $p,q\in\Z$ con $a^{2}=2p+1$ y $b^{2}=2q+1$.
  \item Entonces $a^{2}+b^{2}=2(p+q+1)$ es par, mientras que $c^{2}$ es impar.
  \item Como ningún entero es par e impar a la vez, $a^{2}+b^{2}\neq c^{2}$. $\blacksquare$
\end{enumerate}
\resultado \textbf{en toda terna pitagórica entera hay al menos un número par.}
\comprobacion $(3,4,5)$: $4$ es par; $(5,12,13)$: $12$ es par. No hace falta desplegar «$c^{2}$ es impar»: basta citar la \cref{prop:impar-cuadrado}.
\end{solucion}

\begin{solucion}{eje:3-equivalencia}
\datos $x\equiv_k y\iff\exists q\in\Z\ (x=y+kq)$.
\planteamiento en cada apartado, exhibir el entero $q$ que pide la definición.
\pasos
\begin{enumerate}[label=(\alph*)]
  \item Sea $x\in\Z$. Con $q=0$, $x=x+k\cdot 0$; luego $x\equiv_k x$.
  \item Supóngase $x\equiv_k y$: existe $q\in\Z$ con $x=y+kq$. Entonces $y=x+k(-q)$ y $-q\in\Z$, luego $y\equiv_k x$.
  \item Supóngase $x\equiv_k y$ y $y\equiv_k z$: existen $q,r\in\Z$ con $x=y+kq$ e $y=z+kr$. Sustituyendo, $x=z+kr+kq=z+k(q+r)$ con $q+r\in\Z$, luego $x\equiv_k z$. $\blacksquare$
\end{enumerate}
\resultado \textbf{$\equiv_k$ es reflexiva, simétrica y transitiva} (una relación de equivalencia).
\comprobacion $\equiv_k$ coincide con la congruencia de la \cref{def:congruencia}: $x=y+kq\iff k\divides(x-y)$. Ejemplo: $17\equiv_5 2$ ($q=3$) y $2\equiv_5 17$ ($q=-3$).
\end{solucion}

\begin{solucion}{eje:3-real}
\datos $x,y\in\R$ con $y^{3}+yx^{2}\le x^{3}+xy^{2}$.
\planteamiento contrarrecíproco: si $y>x$, entonces $y^{3}+yx^{2}>x^{3}+xy^{2}$. Factorizar la diferencia.
\pasos
\begin{enumerate}[label=\arabic*.]
  \item Se demuestra el contrarrecíproco. Sean $x,y\in\R$ con $y>x$.
  \item Se calcula $(y^{3}+yx^{2})-(x^{3}+xy^{2})=(y^{3}-x^{3})-xy(y-x)=(y-x)(y^{2}+xy+x^{2})-xy(y-x)=(y-x)(x^{2}+y^{2})$.
  \item Como $y>x$, $y-x>0$. Además, $x$ e $y$ no son ambos nulos (son distintos), luego $x^{2}+y^{2}>0$.
  \item El producto de dos positivos es positivo: $y^{3}+yx^{2}>x^{3}+xy^{2}$, que es la negación de la hipótesis. $\blacksquare$
\end{enumerate}
\resultado \textbf{la hipótesis implica $y\le x$.}
\comprobacion $x=1$, $y=2$: $8+2=10>1+4=5$, y $(y-x)(x^{2}+y^{2})=1\cdot 5=5=10-5$.
\end{solucion}

\begin{solucion}{eje:3-bicond}
\datos $n\in\Z$.
\planteamiento dos implicaciones: ($\Leftarrow$) directa; ($\Rightarrow$) por contrarrecíproco.
\pasos
\begin{enumerate}[label=\arabic*.]
  \item ($\Leftarrow$) Sea $n$ par, $n=2k$ con $k\in\Z$. Entonces $n^{2}=4k^{2}$ y $4\divides n^{2}$.
  \item ($\Rightarrow$) Por contrarrecíproco: sea $n$ impar. Por la \cref{prop:impar-cuadrado}, $n^{2}$ es impar.
  \item Si fuera $4\divides n^{2}$, sería $n^{2}=4j=2(2j)$ par, imposible. Luego $4\nodivide n^{2}$.
  \item Se han probado ambas implicaciones, luego $4\divides n^{2}\iff n$ par. $\blacksquare$
\end{enumerate}
\resultado \textbf{$4\divides n^{2}$ si y solo si $n$ es par.}
\comprobacion $n=6$: $36=4\cdot 9$; $n=5$: $25$ no es múltiplo de~4. Coherente con la \cref{prop:cuadrado-mod4}.
\end{solucion}

\begin{solucion}{eje:3-restos}
\datos $3^{100}\bmod 7$, $7^{2026}\bmod 10$, $(5^{40}+6^{40})\bmod 11$.
\planteamiento encontrar una potencia congruente con $\pm1$ y aplicar la \cref{prop:congruencias}.
\pasos
\begin{enumerate}[label=(\alph*)]
  \item $3^{3}=27\equiv -1\pmod 7$, luego $3^{6}\equiv 1$. Como $100=6\cdot 16+4$, $3^{100}\equiv 3^{4}=81=7\cdot 11+4\equiv 4$.
  \item $7^{2}=49\equiv -1\pmod{10}$, luego $7^{4}\equiv 1$. Como $2026=4\cdot 506+2$, $7^{2026}\equiv 7^{2}\equiv 9\pmod{10}$.
  \item $5^{2}=25\equiv 3$, $5^{5}=5^{4}\cdot 5\equiv 9\cdot 5=45\equiv 1\pmod{11}$, luego $5^{40}=(5^{5})^{8}\equiv 1$. $6^{2}=36\equiv 3$, $6^{5}=6^{4}\cdot 6\equiv 9\cdot 6=54\equiv -1$, luego $6^{40}=(6^{5})^{8}\equiv 1$. La suma es $\equiv 2$.
\end{enumerate}
\resultado (a) \textbf{4}; (b) \textbf{9}; (c) \textbf{2}.
\comprobacion cálculo exacto con Python: \texttt{pow(3,100,7)=4}, \texttt{pow(7,2026,10)=9}, $(5^{40}+6^{40})\bmod 11=2$.
\end{solucion}

\begin{solucion}{eje:3-division}
\datos $2026$ y $7$; $-17$ y $5$; $2023$, $2026$; módulo~9.
\planteamiento buscar $q$ con $qb\le a<(q+1)b$ y tomar $r=a-qb$ (\cref{teo:division}); aplicar la \cref{def:divide}.
\pasos
\begin{enumerate}[label=\arabic*.]
  \item $7\cdot 289=2023\le 2026<2030=7\cdot 290$: $q=289$, $r=2026-2023=3$.
  \item $5\cdot(-4)=-20\le -17<-15=5\cdot(-3)$: $q=-4$, $r=-17+20=3$. (No vale $q=-3$, $r=-2$, porque el resto debe ser no negativo.)
  \item $7\divides 2023$ si existe $c\in\Z$ con $2023=7c$: sí, $c=289$. $7\divides 2026$ exigiría $2026=7c$, pero por la unicidad de la división $2026=7\cdot 289+3$ con resto $3\neq 0$: no.
  \item $2026=9\cdot 225+1$, resto $1$. Suma de cifras: $2+0+2+6=10=9+1$, resto $1$. Coinciden porque $10\equiv 1\pmod 9$, luego $10^{k}\equiv 1$ y cada cifra $d$ en la posición $k$ aporta $d\cdot 10^{k}\equiv d$ (\cref{prop:congruencias}).
\end{enumerate}
\resultado (a) $2026=7\cdot 289+3$; $-17=5\cdot(-4)+3$. (b) $7\divides 2023$, $7\nodivide 2026$. (c) resto $1$ en ambos casos.
\comprobacion Python: \texttt{divmod(2026,7)=(289,3)}, \texttt{divmod(-17,5)=(-4,3)}, \texttt{2026\%9=1}.
\end{solucion}

% ---------------------------------------------------------------------
\section{Tema 4}

\begin{solucion}{eje:4-suponer}
\datos cuatro enunciados.
\planteamiento se supone la negación (\cref{def:absurdo}); para un condicional, hipótesis y negación de la tesis.
\pasos
\begin{enumerate}[label=(\alph*)]
  \item Se supone que $\sqrt5$ es racional: $\sqrt5=a/b$ con $a,b\in\Z$, $b>0$, fracción irreducible.
  \item Se supone que existe $n\in\N$ con $3\divides n^{2}+1$.
  \item Se suponen $a,b\in\Z$ con $ab$ impar y con $a$ par o $b$ par.
  \item Se supone que todo primo es menor o igual que $10^{100}$.
\end{enumerate}
\resultado las cuatro suposiciones anteriores.
\comprobacion en (c), la negación de «$a$ y $b$ impares» es «$a$ par o $b$ par» (De Morgan); en (b), la negación de $\forall n\,\neg Q(n)$ es $\exists n\,Q(n)$.
\end{solucion}

\begin{solucion}{eje:4-raiz3}
\datos $n\in\Z$; se quiere $3\divides n^{2}\Rightarrow 3\divides n$ y $\sqrt3\notin\Q$.
\planteamiento lema por contrarrecíproco con los restos módulo~3; teorema imitando el \cref{teo:raiz2}.
\pasos
\begin{enumerate}[label=\arabic*.]
  \item \emph{Lema.} Se prueba el contrarrecíproco: si $3\nodivide n$, entonces $3\nodivide n^{2}$. Si $3\nodivide n$, por la división euclídea $n=3k+1$ o $n=3k+2$ con $k\in\Z$.
  \item Caso $n=3k+1$: $n^{2}=9k^{2}+6k+1=3(3k^{2}+2k)+1$. Caso $n=3k+2$: $n^{2}=9k^{2}+12k+4=3(3k^{2}+4k+1)+1$. En ambos, $n^{2}$ deja resto~1 al dividir por~3, luego $3\nodivide n^{2}$ (por la unicidad del resto).
  \item \emph{Teorema.} Supóngase, por reducción al absurdo, que $\sqrt3=a/b$ con $a,b\in\Z$, $b>0$ y fracción irreducible. Entonces $a^{2}=3b^{2}$, luego $3\divides a^{2}$ y, por el lema, $3\divides a$: $a=3c$.
  \item Sustituyendo, $9c^{2}=3b^{2}$, es decir, $b^{2}=3c^{2}$, luego $3\divides b^{2}$ y $3\divides b$.
  \item Así $3$ divide a $a$ y a $b$, en contra de que la fracción sea irreducible. Por tanto, $\sqrt3$ es irracional. $\blacksquare$
\end{enumerate}
\resultado \textbf{$3\divides n^{2}\Rightarrow 3\divides n$, y $\sqrt3$ es irracional.}
\comprobacion los cuadrados de $0,1,2$ son $0,1,4$, con restos $0,1,1$ módulo~3: solo los múltiplos de~3 tienen cuadrado múltiplo de~3.
\end{solucion}

\begin{solucion}{eje:4-raiz-cubica}
\datos $\sqrt[3]{2}$.
\planteamiento reducción al absurdo; se necesita «$n^{3}$ par $\Rightarrow n$ par», que es el contrarrecíproco del \cref{eje:3-cubo}.
\pasos
\begin{enumerate}[label=\arabic*.]
  \item \emph{Lema.} Si $n^{3}$ es par, $n$ es par: si $n$ fuera impar, $n^{3}$ sería impar por el \cref{eje:3-cubo}.
  \item Supóngase, por reducción al absurdo, que $\sqrt[3]{2}=a/b$ con $a,b\in\Z$, $b>0$ y fracción irreducible. Elevando al cubo, $a^{3}=2b^{3}$.
  \item Entonces $a^{3}$ es par y, por el lema, $a=2c$ con $c\in\Z$. Sustituyendo, $8c^{3}=2b^{3}$, es decir, $b^{3}=4c^{3}=2(2c^{3})$ es par, y $b$ es par.
  \item Así $2$ divide a $a$ y a $b$, contradicción con la irreducibilidad. Luego $\sqrt[3]{2}$ es irracional. $\blacksquare$
\end{enumerate}
\resultado \textbf{$\sqrt[3]{2}$ es irracional.}
\comprobacion la prueba es la del \cref{teo:raiz2} cambiando cuadrados por cubos; numéricamente $\sqrt[3]2\approx 1{,}2599$ y ninguna fracción de denominador pequeño cumple $a^{3}=2b^{3}$ (por ejemplo, $5^{3}=125\neq 2\cdot 4^{3}=128$).
\end{solucion}

\begin{solucion}{eje:4-log}
\datos $\log_2 3$, el número $x$ con $2^{x}=3$.
\planteamiento suponer $x=a/b$ y obtener una igualdad de enteros con paridades distintas.
\pasos
\begin{enumerate}[label=\arabic*.]
  \item Como $2^{1}<3<2^{2}$, se tiene $1<\log_2 3<2$; en particular, $\log_2 3>0$.
  \item Supóngase, por reducción al absurdo, que $\log_2 3=a/b$ con $a,b\in\N$ (ambos positivos, pues el cociente lo es y se puede tomar $b>0$).
  \item Entonces $2^{a/b}=3$ y, elevando a $b$, $2^{a}=3^{b}$.
  \item Como $a\ge 1$, $2^{a}$ es par; como $3^{b}$ es producto de impares, es impar (\cref{prop:congruencias}(iii) con $3\equiv 1\pmod 2$). Un entero no puede ser par e impar: contradicción.
  \item Por tanto, $\log_2 3$ es irracional. $\blacksquare$
\end{enumerate}
\resultado \textbf{$\log_2 3$ es irracional.}
\comprobacion $\log_2 3\approx 1{,}585$; la hipótesis $a\ge 1$ (paso~2) es esencial, pues $2^{0}=3^{0}$.
\end{solucion}

\begin{solucion}{eje:4-racional-irracional}
\datos $a\in\Q$, $b\notin\Q$.
\planteamiento reducción al absurdo: si $a+b$ (o $ab$) fuera racional, $b$ se despejaría como racional.
\pasos
\begin{enumerate}[label=\arabic*.]
  \item Escríbase $a=m/n$ con $m,n\in\Z$, $n\neq 0$.
  \item \emph{Suma.} Supóngase que $a+b=p/q$ con $p,q\in\Z$, $q\neq 0$. Entonces $b=\frac{p}{q}-\frac{m}{n}=\frac{pn-mq}{qn}$, con $pn-mq\in\Z$ y $qn\neq 0$; luego $b$ es racional, contradicción.
  \item \emph{Producto}, con $a\neq 0$ (es decir, $m\neq 0$). Supóngase $ab=p/q$. Entonces $b=\frac{p}{q}\cdot\frac{n}{m}=\frac{pn}{qm}$ con $qm\neq 0$, luego $b$ es racional, contradicción.
  \item Si $a=0$, $ab=0$ es racional: la hipótesis $a\neq 0$ es necesaria. $\blacksquare$
\end{enumerate}
\resultado \textbf{$a+b$ es irracional; $ab$ es irracional si $a\neq 0$}, y racional ($=0$) si $a=0$.
\comprobacion ejemplos: $1+\sqrt2$ y $3\sqrt2$ son irracionales; $0\cdot\sqrt2=0$.
\end{solucion}

\begin{solucion}{eje:4-diofantica}
\datos $21a+30b=1$ y $21a+30b=3$.
\planteamiento $3$ divide a $21$ y a $30$, luego a toda combinación $21a+30b$ (\cref{prop:divisibilidad}(ii)).
\pasos
\begin{enumerate}[label=\arabic*.]
  \item Supóngase, por reducción al absurdo, que existen $a,b\in\Z$ con $21a+30b=1$.
  \item Como $21=3\cdot 7$ y $30=3\cdot 10$, $21a+30b=3(7a+10b)$, luego $3\divides 1$.
  \item Pero los únicos divisores de~1 son $\pm1$, y $3\neq\pm1$. Contradicción; no existen tales enteros. $\blacksquare$
  \item Para el valor~3 sí hay solución: $a=3$, $b=-2$, pues $63-60=3$.
\end{enumerate}
\resultado \textbf{no hay solución para~1; sí para~3} (por ejemplo, $(a,b)=(3,-2)$).
\comprobacion $\operatorname{mcd}(21,30)=3$: las combinaciones enteras de $21$ y $30$ son exactamente los múltiplos de~3 (identidad de Bézout, que excede este módulo).
\end{solucion}

\begin{solucion}{eje:4-mod4}
\datos $n\in\Z$; enteros $4k+3$.
\planteamiento \cref{prop:cuadrado-mod4}: todo cuadrado es $\equiv 0$ o $1\pmod 4$.
\pasos
\begin{enumerate}[label=\arabic*.]
  \item Sea $n\in\Z$. Por la \cref{prop:cuadrado-mod4}, $n^{2}\equiv 0$ o $n^{2}\equiv 1\pmod 4$.
  \item Caso $n^{2}\equiv 0$: $n^{2}+2\equiv 2\pmod 4$. Caso $n^{2}\equiv 1$: $n^{2}+2\equiv 3\pmod 4$. En ningún caso es $\equiv 0$, luego $4\nodivide n^{2}+2$.
  \item Sean ahora $x,y\in\Z$. Por la misma proposición, $x^{2}+y^{2}$ es congruente con $0+0$, $0+1$, $1+0$ o $1+1$, es decir, con $0$, $1$ o $2\pmod 4$.
  \item Un entero $4k+3$ es $\equiv 3\pmod 4$, que no está entre esos valores. Por la unicidad del resto, $x^{2}+y^{2}\neq 4k+3$. $\blacksquare$
\end{enumerate}
\resultado \textbf{$4\nodivide n^{2}+2$ y ningún entero $4k+3$ es suma de dos cuadrados.}
\comprobacion $3$, $7$, $11$, $15$ no son sumas de dos cuadrados; $5=1+4$ y $13=4+9$ sí, y son $\equiv 1\pmod 4$.
\end{solucion}

\begin{solucion}{eje:4-24}
\datos $p$ primo, $p>3$.
\planteamiento $p^{2}-1=(p-1)(p+1)$; probar por separado $8\divides p^{2}-1$ y $3\divides p^{2}-1$.
\pasos
\begin{enumerate}[label=\arabic*.]
  \item Sea $p>3$ primo. No es par (sería divisible por~2 y distinto de~2) ni múltiplo de~3 (sería divisible por~3 y distinto de~3).
  \item \emph{Factor 8.} Como $p$ es impar, $p=2k+1$ y $p^{2}-1=(p-1)(p+1)=2k(2k+2)=4k(k+1)$. Por la \cref{prop:consecutivos}, $k(k+1)=k^{2}+k$ es par, $k(k+1)=2j$, luego $p^{2}-1=8j$.
  \item \emph{Factor 3.} Como $3\nodivide p$, $p=3m+1$ o $p=3m+2$. En el primer caso, $3\divides p-1$; en el segundo, $p+1=3(m+1)$ y $3\divides p+1$. En ambos, $3\divides(p-1)(p+1)=p^{2}-1$.
  \item Por los pasos 2 y 3, $3\divides 8j$. Como también $3\divides 9j$, la \cref{prop:divisibilidad}(ii) da $3\divides 9j-8j=j$. Así $j=3i$ con $i\in\Z$ y $p^{2}-1=24i$. $\blacksquare$
\end{enumerate}
\resultado \textbf{$24\divides p^{2}-1$ para todo primo $p>3$.}
\comprobacion $5^{2}-1=24$, $7^{2}-1=48$, $11^{2}-1=120=24\cdot 5$; se ha comprobado con Python para todos los primos entre $5$ y $10^{4}$. Para $p=3$ falla ($8$), lo que muestra que la hipótesis $p>3$ se usa (paso~3).
\end{solucion}

\begin{solucion}{eje:4-maximo}
\datos $x,y\in\R$.
\planteamiento dos casos según el signo de $x-y$.
\pasos
\begin{enumerate}[label=\arabic*.]
  \item Caso $x\ge y$: $\abs{x-y}=x-y$, luego $\frac{x+y+\abs{x-y}}{2}=\frac{2x}{2}=x=\max(x,y)$ y $\frac{x+y-\abs{x-y}}{2}=y=\min(x,y)$.
  \item Caso $x<y$: $\abs{x-y}=y-x$, luego la primera expresión vale $\frac{2y}{2}=y=\max(x,y)$ y la segunda $x=\min(x,y)$.
  \item Comprobación con $(3,-5)$: $x+y=-2$, $\abs{x-y}=8$; $\max=(-2+8)/2=3$, $\min=(-2-8)/2=-5$.
\end{enumerate}
\resultado \textbf{ambas fórmulas son ciertas}; para $(3,-5)$ dan $3$ y $-5$.
\comprobacion sumando las fórmulas, $\max+\min=x+y$; restándolas, $\max-\min=\abs{x-y}$, ambas evidentes.
\end{solucion}

\begin{solucion}{eje:4-restos-cuadrados}
\datos $n\in\Z$; módulos $3$, $5$ y $8$.
\planteamiento por la \cref{prop:congruencias}, $n^{2}\bmod m$ solo depende del resto de $n$ módulo $m$; se recorren todos los restos.
\pasos
\begin{enumerate}[label=\arabic*.]
  \item Módulo 3, restos $0,1,2$: cuadrados $0,1,4\equiv 1$. Valores posibles: $\{0,1\}$.
  \item Módulo 5, restos $0,\dots,4$: cuadrados $0,1,4,9\equiv 4,16\equiv 1$. Valores posibles: $\{0,1,4\}$.
  \item Módulo 8, restos $0,\dots,7$: cuadrados $0,1,4,9\equiv1,16\equiv0,25\equiv1,36\equiv4,49\equiv1$. Valores: $\{0,1,4\}$; los restos impares $1,3,5,7$ dan siempre~$1$.
  \item (a) $n^{2}+1\equiv 1$ o $2\pmod 3$, nunca $0$. (b) $2$ y $3$ no están en $\{0,1,4\}$. (c) Un impar tiene resto impar módulo~8 y, por el paso~3, $n^{2}\equiv 1$.
\end{enumerate}
\resultado cuadrados: módulo~3 $\{0,1\}$; módulo~5 $\{0,1,4\}$; módulo~8 $\{0,1,4\}$. Las tres afirmaciones son ciertas.
\comprobacion $7^{2}=49=8\cdot 6+1$; $13^{2}=169=5\cdot 33+4$; $5^{2}+1=26$, no múltiplo de~3. El apartado (c) reaparece en el \cref{eje:4-24}.
\end{solucion}

\begin{solucion}{eje:4-triangular-inversa}
\datos $x,y\in\R$.
\planteamiento aplicar el \cref{teo:triangular} dos veces y combinar con $\abs{t}\le c\iff -c\le t\le c$.
\pasos
\begin{enumerate}[label=\arabic*.]
  \item Sean $x,y\in\R$. Como $x=(x-y)+y$, la desigualdad triangular da $\abs{x}\le\abs{x-y}+\abs{y}$, es decir, $\abs{x}-\abs{y}\le\abs{x-y}$.
  \item Como $y=(y-x)+x$, del mismo modo $\abs{y}-\abs{x}\le\abs{y-x}=\abs{x-y}$, es decir, $-\abs{x-y}\le\abs{x}-\abs{y}$.
  \item Juntando, $-\abs{x-y}\le\abs{x}-\abs{y}\le\abs{x-y}$. Si $t=\abs{x}-\abs{y}\ge 0$, $\abs{t}=t\le\abs{x-y}$; si $t<0$, $\abs{t}=-t\le\abs{x-y}$. En ambos casos, $\bigl\lvert\abs{x}-\abs{y}\bigr\rvert\le\abs{x-y}$. $\blacksquare$
\end{enumerate}
\resultado \textbf{desigualdad triangular inversa demostrada.}
\comprobacion $x=3$, $y=-5$: $\bigl\lvert 3-5\bigr\rvert=2\le\abs{8}=8$; $x=3$, $y=5$: $2\le 2$ (igualdad cuando tienen el mismo signo).
\end{solucion}

\begin{solucion}{eje:4-producto-par}
\datos $m,n\in\Z$, $mn$ par, $m$ impar.
\planteamiento la negación de $P\Rightarrow Q$ es $P\wedge\neg Q$: suponer $mn$ par, $m$ impar y $n$ impar.
\pasos
\begin{enumerate}[label=\arabic*.]
  \item Supóngase, por reducción al absurdo, que existen enteros $m$ y $n$ con $mn$ par, $m$ impar y $n$ no par, es decir, impar.
  \item Existen $k,r\in\Z$ con $m=2k+1$ y $n=2r+1$.
  \item Entonces $mn=4kr+2k+2r+1=2(2kr+k+r)+1$ es impar.
  \item Esto contradice que $mn$ sea par, pues ningún entero es par e impar. Luego, si $mn$ es par y $m$ impar, $n$ es par. $\blacksquare$
\end{enumerate}
\resultado \textbf{$n$ es par.}
\comprobacion $m=3$, $mn=12$: $n=4$ es par. La prueba usa la suposición «$n$ impar» en el paso~2, de modo que no es un «sándwich de contradicción».
\end{solucion}

\begin{solucion}{eje:4-cubos}
\datos todo entero $x$; se quieren cubos de enteros.
\planteamiento (a) identidad $(n+1)^{3}+(n-1)^{3}=2n^{3}+6n$; (b) $x^{3}-x=(x-1)x(x+1)$; (c) combinar.
\pasos
\begin{enumerate}[label=\arabic*.]
  \item (a) Sea $N=6n$ con $n\in\Z$. Se desarrolla $(n+1)^{3}+(n-1)^{3}=(n^{3}+3n^{2}+3n+1)+(n^{3}-3n^{2}+3n-1)=2n^{3}+6n$.
  \item Por tanto, $6n=(n+1)^{3}+(n-1)^{3}+(-n)^{3}+(-n)^{3}$, suma de cuatro cubos de enteros.
  \item (b) Sea $x\in\Z$. $x^{3}-x=(x-1)x(x+1)$ es producto de tres enteros consecutivos. Entre $x-1$ y $x$ hay uno par, de modo que $2\divides x^{3}-x$. Por casos según el resto de $x$ módulo~3: si $x\equiv 0$, $3\divides x$; si $x\equiv 1$, $3\divides x-1$; si $x\equiv 2$, $3\divides x+1$. Luego $3\divides x^{3}-x$.
  \item Escríbase $x^{3}-x=2s$ con $s\in\Z$. Como $3\divides 2s$ y $3\divides 3s$, se tiene $3\divides 3s-2s=s$ (\cref{prop:divisibilidad}(ii)); así $s=3t$ y $x^{3}-x=6t$, es decir, $x^{3}\equiv x\pmod 6$.
  \item (c) Sea $x\in\Z$. Por (b), $x^{3}-x=6t$ con $t\in\Z$, luego $x=x^{3}+6(-t)$. Por (a), $6(-t)$ es suma de cuatro cubos, y $x$ es suma de $x^{3}$ y esos cuatro: cinco cubos. $\blacksquare$
  \item Para $x=7$: $7^{3}-7=336=6\cdot 56$, $t=56$, $-t=-56$, y $6(-56)=(-55)^{3}+(-57)^{3}+56^{3}+56^{3}$. Así, $7=7^{3}+(-55)^{3}+(-57)^{3}+56^{3}+56^{3}$.
\end{enumerate}
\resultado \textbf{todo entero es suma de cinco cubos}; $7=7^{3}+(-55)^{3}+(-57)^{3}+2\cdot 56^{3}$.
\comprobacion $343-166\,375-185\,193+175\,616+175\,616=7$ (verificado con Python). No se sabe si basta con cuatro cubos para todo entero~\cite{cs103}.
\end{solucion}

\begin{solucion}{eje:4-euclides-n}
\datos los seis primeros primos $2,3,5,7,11,13$.
\planteamiento calcular los productos y comprobar la primalidad dividiendo por los primos hasta la raíz cuadrada.
\pasos
\begin{enumerate}[label=\arabic*.]
  \item $N_1=3$, $N_2=7$, $N_3=31$, $N_4=211$, $N_5=2311$, $N_6=30\,031$.
  \item $3,7,31$ son primos. $211$: no es divisible por $2,3,5,7,11,13$ y $\sqrt{211}<15$: primo. $2311$: no es divisible por ningún primo hasta $47$ y $\sqrt{2311}<49$: primo.
  \item $30\,031=59\cdot 509$: compuesto.
  \item No hay contradicción: la demostración solo afirma que los divisores primos de $N_k$ no están entre $p_1,\dots,p_k$, y en efecto $59$ y $509$ son mayores que~13.
\end{enumerate}
\resultado $N_1,\dots,N_5$ \textbf{primos}; $N_6=\mathbf{59\cdot 509}$.
\comprobacion $59\cdot 509=30\,031$; factorización confirmada con Python (\texttt{sympy.factorint}).
\end{solucion}

\begin{solucion}{eje:4-4k3}
\datos primos $\equiv 3\pmod 4$; hay al menos uno, el~3.
\planteamiento reducción al absurdo con $N=4p_1\cdots p_m-1$.
\pasos
\begin{enumerate}[label=\arabic*.]
  \item Supóngase que solo hay un número finito de primos de la forma $4k+3$, y sean $p_1,\dots,p_m$ todos ellos ($m\ge 1$, pues $3$ es uno). Sea $N=4p_1\cdots p_m-1$. Entonces $N\ge 4\cdot 3-1=11$ y $N\equiv -1\equiv 3\pmod 4$; en particular, $N$ es impar.
  \item Por el \cref{teo:factorizacion}, $N$ es producto de primos $q_1\cdots q_s$ (con $s=1$ si $N$ es primo). Todos son impares porque $N$ lo es, luego cada $q_j\equiv 1$ o $q_j\equiv 3\pmod 4$.
  \item Si todos fueran $\equiv 1\pmod 4$, su producto sería $\equiv 1\cdot1\cdots1=1\pmod 4$ (\cref{prop:congruencias}(ii)), en contra de $N\equiv 3$. Luego algún $q_j\equiv 3\pmod 4$.
  \item Por la hipótesis, $q_j=p_i$ para algún $i$. Entonces $q_j$ divide a $N$ y a $4p_1\cdots p_m$, luego divide a su diferencia, $1$. Imposible, porque $q_j\ge 3$.
  \item Por tanto, hay infinitos primos de la forma $4k+3$. $\blacksquare$
\end{enumerate}
\resultado \textbf{existen infinitos primos $\equiv 3\pmod 4$.}
\comprobacion con la lista $\{3\}$: $N=11$, primo y $\equiv 3$; con $\{3,7,11\}$: $N=923=13\cdot 71$, y $71\equiv 3\pmod 4$ no está en la lista. (La analogía con $4p_1\cdots p_m+1$ no funciona para los primos $4k+1$: un producto de primos $\equiv 3$ puede ser $\equiv 1$.)
\end{solucion}

\begin{solucion}{eje:4-raiz-primo}
\datos $p$ primo; lema de Euclides.
\planteamiento como en el \cref{teo:raiz2}; el lema de Euclides con $a=b=n$ da $p\divides n^{2}\Rightarrow p\divides n$.
\pasos
\begin{enumerate}[label=\arabic*.]
  \item Supóngase, por reducción al absurdo, que $\sqrt p=a/b$ con $a,b\in\Z$, $b>0$ y fracción irreducible. Entonces $a^{2}=pb^{2}$, luego $p\divides a\cdot a$.
  \item Por el lema de Euclides, $p\divides a$: $a=pc$ con $c\in\Z$. Sustituyendo, $p^{2}c^{2}=pb^{2}$, es decir, $b^{2}=pc^{2}$, y de nuevo $p\divides b$.
  \item Así $p$ divide a $a$ y a $b$, y como $p\ge 2$, contradice la irreducibilidad. Luego $\sqrt p$ es irracional. $\blacksquare$
  \item Con $4$ (no primo) falla el paso~2: $4\divides 2\cdot 2$ pero $4\nodivide 2$. De hecho $\sqrt4=2/1$ es racional.
\end{enumerate}
\resultado \textbf{$\sqrt p$ es irracional para todo primo $p$}; con $\sqrt4$ falla el lema de Euclides.
\comprobacion los casos $p=2$ y $p=3$ son el \cref{teo:raiz2} y el \cref{eje:4-raiz3}, donde el lema de Euclides se sustituyó por un argumento de restos.
\end{solucion}

% ---------------------------------------------------------------------
\section{Tema 5}

\begin{solucion}{eje:5-plantilla}
\datos $\sum_{k=1}^{n}(2k-1)=n^{2}$ para $n\ge 1$.
\planteamiento identificar las piezas del \cref{teo:induccion} con caso base $n_0=1$.
\pasos
\begin{enumerate}[label=\arabic*.]
  \item (a) $P(n)$: «$1+3+\dots+(2n-1)=n^{2}$».
  \item (b) $P(1)$: «$1=1^{2}$». $P(n+1)$: «$1+3+\dots+(2n-1)+(2n+1)=(n+1)^{2}$».
  \item (c) Hay que demostrar: para todo $n\ge 1$, si $P(n)$ es cierta, también lo es $P(n+1)$.
  \item Demostración. \emph{Caso base:} $P(1)$ es $1=1$. \emph{Paso:} sea $n\ge 1$ con $P(n)$. Entonces $1+\dots+(2n-1)+(2n+1)=n^{2}+2n+1=(n+1)^{2}$, que es $P(n+1)$. Por inducción, $P(n)$ vale para todo $n\ge 1$. $\blacksquare$
\end{enumerate}
\resultado \textbf{la suma de los $n$ primeros impares es $n^{2}$.}
\comprobacion $1+3+5+7=16=4^{2}$.
\end{solucion}

\begin{solucion}{eje:5-cuadrados}
\datos $P(n)$: $\sum_{k=1}^{n}k^{2}=\frac{n(n+1)(2n+1)}{6}$.
\planteamiento inducción simple desde $n=1$.
\pasos
\begin{enumerate}[label=\arabic*.]
  \item \emph{Caso base.} $n=1$: $1^{2}=1$ y $\frac{1\cdot 2\cdot 3}{6}=1$.
  \item \emph{Paso.} Sea $n\ge 1$ y supóngase $P(n)$. Entonces
  \[
  \begin{aligned}
  \sum_{k=1}^{n+1}k^{2}&=\frac{n(n+1)(2n+1)}{6}+(n+1)^{2}\\
  &=\frac{(n+1)\bigl[n(2n+1)+6(n+1)\bigr]}{6}=\frac{(n+1)(2n^{2}+7n+6)}{6}.
  \end{aligned}
  \]
  \item Como $2n^{2}+7n+6=(n+2)(2n+3)$, la suma es $\frac{(n+1)(n+2)(2(n+1)+1)}{6}$, que es $P(n+1)$.
  \item Por el principio de inducción, $P(n)$ vale para todo $n\in\N$. $\blacksquare$
\end{enumerate}
\resultado $\sum_{k=1}^{n}k^{2}=\dfrac{n(n+1)(2n+1)}{6}$.
\comprobacion $n=4$: $1+4+9+16=30$ y $\frac{4\cdot5\cdot9}{6}=30$. Factorización: $(n+2)(2n+3)=2n^{2}+3n+4n+6$.
\end{solucion}

\begin{solucion}{eje:5-cubos}
\datos $P(n)$: $\sum_{k=1}^{n}k^{3}=\bigl(\frac{n(n+1)}{2}\bigr)^{2}$.
\planteamiento inducción simple; sacar factor común $(n+1)^{2}/4$.
\pasos
\begin{enumerate}[label=\arabic*.]
  \item \emph{Caso base.} $n=1$: $1=(1\cdot 2/2)^{2}=1$.
  \item \emph{Paso.} Sea $n\ge 1$ con $P(n)$. Entonces
  \[
  \sum_{k=1}^{n+1}k^{3}=\frac{n^{2}(n+1)^{2}}{4}+(n+1)^{3}=\frac{(n+1)^{2}\bigl[n^{2}+4(n+1)\bigr]}{4}=\frac{(n+1)^{2}(n+2)^{2}}{4},
  \]
  que es $P(n+1)$.
  \item Por inducción, $P(n)$ vale para todo $n\in\N$. $\blacksquare$
\end{enumerate}
\resultado $\sum_{k=1}^{n}k^{3}=\Bigl(\dfrac{n(n+1)}{2}\Bigr)^{2}$, el cuadrado de la suma del \cref{ej:gauss}.
\comprobacion $n=3$: $1+8+27=36=6^{2}$.
\end{solucion}

\begin{solucion}{eje:5-telescopica}
\datos $P(n)$: $\sum_{k=1}^{n}\frac{1}{k(k+1)}=\frac{n}{n+1}$.
\planteamiento inducción; después, suma telescópica.
\pasos
\begin{enumerate}[label=\arabic*.]
  \item \emph{Caso base.} $n=1$: $\frac{1}{2}=\frac{1}{2}$.
  \item \emph{Paso.} Si vale $P(n)$, entonces $\sum_{k=1}^{n+1}\frac{1}{k(k+1)}=\frac{n}{n+1}+\frac{1}{(n+1)(n+2)}=\frac{n(n+2)+1}{(n+1)(n+2)}=\frac{(n+1)^{2}}{(n+1)(n+2)}=\frac{n+1}{n+2}$, que es $P(n+1)$. Por inducción, vale para todo $n$. $\blacksquare$
  \item \emph{Segunda prueba.} Como $\frac1k-\frac1{k+1}=\frac{k+1-k}{k(k+1)}=\frac{1}{k(k+1)}$, la suma es $\bigl(1-\tfrac12\bigr)+\bigl(\tfrac12-\tfrac13\bigr)+\dots+\bigl(\tfrac1n-\tfrac1{n+1}\bigr)$; cada término intermedio aparece una vez con signo $+$ y otra con signo $-$, y queda $1-\frac{1}{n+1}=\frac{n}{n+1}$.
\end{enumerate}
\resultado $\sum_{k=1}^{n}\dfrac{1}{k(k+1)}=\dfrac{n}{n+1}$.
\comprobacion $n=3$: $\frac12+\frac16+\frac1{12}=\frac{6+2+1}{12}=\frac{9}{12}=\frac34$. El argumento telescópico, escrito con rigor, es a su vez una inducción.
\end{solucion}

\begin{solucion}{eje:5-potencia}
\datos $P(n)$: $2^{n}>n^{2}$, para $n\ge 5$.
\planteamiento inducción con caso base $n=5$; en el paso hay que probar $2n^{2}\ge(n+1)^{2}$ para $n\ge 5$.
\pasos
\begin{enumerate}[label=\arabic*.]
  \item \emph{Caso base.} $2^{5}=32>25=5^{2}$.
  \item \emph{Desigualdad auxiliar.} Para $n\ge 5$: $2n^{2}-(n+1)^{2}=n^{2}-2n-1=(n-1)^{2}-2\ge 4^{2}-2>0$.
  \item \emph{Paso.} Sea $n\ge 5$ con $2^{n}>n^{2}$. Entonces $2^{n+1}=2\cdot 2^{n}>2n^{2}>(n+1)^{2}$, por el paso~2.
  \item Por inducción, $2^{n}>n^{2}$ para todo $n\ge 5$. $\blacksquare$
  \item Valores pequeños: $n=1$: $2>1$ (cierta); $n=2$: $4>4$ (falsa); $n=3$: $8>9$ (falsa); $n=4$: $16>16$ (falsa).
\end{enumerate}
\resultado \textbf{cierta para $n\ge 5$ y para $n=1$; falsa para $n=2,3,4$.}
\comprobacion $n=6$: $64>36$; $n=10$: $1024>100$. El paso auxiliar falla para $n=2$ ($(n-1)^{2}-2=-1$), lo que explica que no pueda empezarse antes.
\end{solucion}

\begin{solucion}{eje:5-div9}
\datos $P(n)$: $9\divides 4^{n}+6n-1$.
\planteamiento inducción; expresar $4^{n+1}+6(n+1)-1$ en función de $4^{n}+6n-1$.
\pasos
\begin{enumerate}[label=\arabic*.]
  \item \emph{Caso base.} $n=0$: $1+0-1=0=9\cdot 0$.
  \item \emph{Paso.} Sea $n\ge 0$ con $4^{n}+6n-1=9m$, $m\in\Z$. Entonces $4^{n}=9m-6n+1$ y
  \[
  4^{n+1}+6(n+1)-1=4(9m-6n+1)+6n+5=36m-18n+9=9(4m-2n+1).
  \]
  \item Como $4m-2n+1\in\Z$, vale $P(n+1)$. Por inducción, $P(n)$ para todo $n\in\N_0$. $\blacksquare$
\end{enumerate}
\resultado \textbf{$9\divides 4^{n}+6n-1$ para todo $n\ge 0$.}
\comprobacion $n=1$: $4+6-1=9$; $n=2$: $16+12-1=27$; $n=3$: $64+18-1=81$. Verificado con Python hasta $n=49$.
\end{solucion}

\begin{solucion}{eje:5-bernoulli}
\datos $x\ge -1$; $P(n)$: $(1+x)^{n}\ge 1+nx$.
\planteamiento inducción; multiplicar la hipótesis por $1+x\ge 0$.
\pasos
\begin{enumerate}[label=\arabic*.]
  \item \emph{Caso base.} $n=0$: $1\ge 1$.
  \item \emph{Paso.} Sea $n\ge 0$ con $(1+x)^{n}\ge 1+nx$. Como $x\ge -1$, el número $1+x$ es no negativo, y multiplicar por él conserva la desigualdad: $(1+x)^{n+1}\ge(1+nx)(1+x)$.
  \item $(1+nx)(1+x)=1+(n+1)x+nx^{2}\ge 1+(n+1)x$, porque $nx^{2}\ge 0$.
  \item Luego $(1+x)^{n+1}\ge 1+(n+1)x$. Por inducción, la desigualdad vale para todo $n\in\N_0$. $\blacksquare$
\end{enumerate}
\resultado \textbf{$(1+x)^{n}\ge 1+nx$ para $x\ge -1$}; la hipótesis se usa en el paso~2.
\comprobacion sin la hipótesis puede fallar: con $x=-4$ y $n=3$, $(-3)^{3}=-27<1-12=-11$. Con $x=0{,}1$ y $n=10$: $1{,}1^{10}\approx 2{,}594\ge 2$.
\end{solucion}

\begin{solucion}{eje:5-sellos}
\datos $n\ge 12$; $n=4a+5b$ con $a,b\ge 0$.
\planteamiento inducción fuerte; el paso añade un~4 al caso $n-3$, que exige cuatro casos base.
\pasos
\begin{enumerate}[label=\arabic*.]
  \item \emph{Casos base.} $12=4\cdot 3$; $13=4\cdot 2+5$; $14=4+5\cdot 2$; $15=5\cdot 3$.
  \item \emph{Paso.} Sea $n\ge 15$ y supóngase el enunciado cierto para $12,13,\dots,n$. Como $12\le n-3\le n$, existen $a,b\in\N_0$ con $n-3=4a+5b$, y entonces $n+1=4(a+1)+5b$.
  \item Por inducción fuerte, todo $n\ge 12$ admite la forma. $\blacksquare$
  \item Menores que~12: los representables son $4,5,8,9,10$; \textbf{no} lo son $1,2,3,6,7,11$.
\end{enumerate}
\resultado \textbf{todo $n\ge 12$ es $4a+5b$}; fallan $1,2,3,6,7,11$.
\comprobacion $11$ no es $4a+5b$: con $b=0,1,2$ quedan $11,6,1$, ninguno múltiplo de~4. Es el mayor no representable, $4\cdot 5-4-5=11$ (fórmula de Sylvester, ampliación).
\end{solucion}

\begin{solucion}{eje:5-fib-suma}
\datos $F_1=F_2=1$, $F_{n+2}=F_{n+1}+F_n$.
\planteamiento inducción simple usando la recurrencia en el paso.
\pasos
\begin{enumerate}[label=\arabic*.]
  \item \emph{Caso base.} $n=1$: $F_1=1$ y $F_3-1=2-1=1$.
  \item \emph{Paso.} Sea $n\ge 1$ con $\sum_{k=1}^{n}F_k=F_{n+2}-1$. Entonces $\sum_{k=1}^{n+1}F_k=F_{n+2}-1+F_{n+1}=F_{n+3}-1$, por la recurrencia.
  \item Por inducción, la igualdad vale para todo $n\in\N$. $\blacksquare$
\end{enumerate}
\resultado $\sum_{k=1}^{n}F_k=F_{n+2}-1$.
\comprobacion $1+1+2+3+5=12=F_7-1=13-1$.
\end{solucion}

\begin{solucion}{eje:5-fib-cota}
\datos $F_n\le(5/3)^{n}$.
\planteamiento inducción fuerte: $F_{n+1}$ depende de $F_n$ y $F_{n-1}$.
\pasos
\begin{enumerate}[label=\arabic*.]
  \item \emph{Casos base.} $F_1=1\le 5/3$ y $F_2=1\le 25/9$.
  \item \emph{Paso.} Sea $n\ge 2$ y supóngase la cota para $1,\dots,n$. Entonces
  \[
  F_{n+1}=F_n+F_{n-1}\le\Bigl(\tfrac53\Bigr)^{n}+\Bigl(\tfrac53\Bigr)^{n-1}=\Bigl(\tfrac53\Bigr)^{n-1}\Bigl(\tfrac53+1\Bigr)=\Bigl(\tfrac53\Bigr)^{n-1}\cdot\tfrac83.
  \]
  \item Como $\frac83=\frac{24}{9}<\frac{25}{9}=\bigl(\frac53\bigr)^{2}$, se obtiene $F_{n+1}\le\bigl(\frac53\bigr)^{n+1}$.
  \item Por inducción fuerte, $F_n\le(5/3)^{n}$ para todo $n\in\N$. $\blacksquare$
  \item Hacen falta dos casos base porque el paso usa $F_{n-1}$: con $n=1$ el paso hablaría de $F_0$, que no está definido en el enunciado. El paso empieza en $n=2$ y produce $F_3$; $F_1$ y $F_2$ deben comprobarse aparte.
\end{enumerate}
\resultado \textbf{$F_n\le(5/3)^{n}$ para todo $n\ge 1$.}
\comprobacion $F_{10}=55\le(5/3)^{10}\approx 165{,}4$. Verificado con Python hasta $n=39$.
\end{solucion}

\begin{solucion}{eje:5-binario}
\datos $n\in\N$.
\planteamiento inducción fuerte; si $n$ es par, $n=2m$ con $m<n$; si es impar mayor que~1, $n=2m+1$ con $m<n$.
\pasos
\begin{enumerate}[label=\arabic*.]
  \item Sea $P(n)$: «$n$ es suma de potencias de~2 distintas». \emph{Caso base:} $1=2^{0}$.
  \item \emph{Paso.} Sea $n\ge 1$ y supóngase $P(1),\dots,P(n)$. Se prueba $P(n+1)$ distinguiendo dos casos.
  \item \emph{Caso $n+1$ par.} $n+1=2m$ con $1\le m\le n$. Por hipótesis, $m=2^{e_1}+\dots+2^{e_r}$ con exponentes distintos. Entonces $n+1=2^{e_1+1}+\dots+2^{e_r+1}$, y los exponentes siguen siendo distintos.
  \item \emph{Caso $n+1$ impar.} Como $n+1\ge 2$ es impar, $n+1\ge 3$ y $n+1=2m+1$ con $1\le m\le n$. Escribiendo $2m=2^{e_1+1}+\dots+2^{e_r+1}$ como antes, todos los exponentes son al menos~1, luego $n+1=2^{0}+2^{e_1+1}+\dots+2^{e_r+1}$ con exponentes distintos.
  \item Por inducción fuerte, $P(n)$ vale para todo $n\in\N$. $\blacksquare$
\end{enumerate}
\resultado \textbf{todo natural tiene representación binaria} (la unicidad también es cierta, aunque no se pedía).
\comprobacion $13$ es impar: $13=2\cdot 6+1$, $6=2\cdot 3$, $3=2\cdot 1+1$; resulta $13=2^{0}+2^{2}+2^{3}$, es decir, $1101_2$. En el módulo~12 esta representación es la base de la codificación de bytes.
\end{solucion}

\begin{solucion}{eje:5-bo-geometrica}
\datos $P(n)$: $\sum_{k=1}^{n}(2k-1)=n^{2}$.
\planteamiento plantilla del menor contraejemplo~\cite[§2.2]{mcs}.
\pasos
\begin{enumerate}[label=\arabic*.]
  \item Sea $C=\{n\in\N:\ P(n)\text{ es falsa}\}$. Supóngase, por reducción al absurdo, que $C\neq\varnothing$.
  \item Por el principio del buen orden, $C$ tiene un elemento mínimo $m$.
  \item $m\neq 1$, porque $P(1)$ es $1=1^{2}$, cierta. Luego $m\ge 2$ y $m-1\in\N$; por la minimalidad de $m$, $m-1\notin C$, es decir, $P(m-1)$ es cierta: $\sum_{k=1}^{m-1}(2k-1)=(m-1)^{2}$.
  \item Sumando $2m-1$: $\sum_{k=1}^{m}(2k-1)=(m-1)^{2}+2m-1=m^{2}$, es decir, $P(m)$ es cierta, en contra de $m\in C$.
  \item Por tanto $C=\varnothing$ y $P(n)$ vale para todo $n\in\N$. $\blacksquare$
\end{enumerate}
\resultado \textbf{$\sum_{k=1}^{n}(2k-1)=n^{2}$}, demostrado por el buen orden.
\comprobacion los pasos 3 y 4 son exactamente el caso base y el paso de inducción del \cref{eje:5-plantilla}, reorganizados: es la idea de la \cref{prop:bo-induccion}.
\end{solucion}

\begin{solucion}{eje:5-caballos}
\datos la «demostración» propuesta.
\planteamiento comprobar el paso de inducción para cada $n$, en particular el menor.
\pasos
\begin{enumerate}[label=\arabic*.]
  \item El paso afirma que los dos grupos de $n$ caballos (sin el primero y sin el último) comparten algún caballo, y que ese caballo común transmite el color de un grupo al otro.
  \item Para $n=1$, el conjunto tiene $n+1=2$ caballos $\{c_1,c_2\}$. Los grupos son $\{c_2\}$ y $\{c_1\}$, que \textbf{no} comparten ningún caballo. No hay forma de relacionar el color de $c_1$ con el de $c_2$.
  \item Para $n\ge 2$ el argumento sí funciona (los grupos comparten $c_2,\dots,c_n$), pero la cadena se rompe en $P(1)\Rightarrow P(2)$, y sin ese eslabón no se deduce nada para $n\ge 2$.
\end{enumerate}
\resultado \textbf{el paso falla para $n=1$} (paso de 1 a 2 caballos), porque los dos grupos son disjuntos.
\comprobacion el enunciado $P(2)$ es falso (dos caballos de colores distintos), coherente con que el paso $P(1)\Rightarrow P(2)$ sea el defectuoso, ya que $P(1)$ es verdadera.
\end{solucion}

\begin{solucion}{eje:5-factorial}
\datos $S_n=\sum_{k=1}^{n}k\cdot k!$.
\planteamiento calcular, conjeturar y demostrar por inducción.
\pasos
\begin{enumerate}[label=\arabic*.]
  \item $S_1=1$; $S_2=1+4=5$; $S_3=5+18=23$; $S_4=23+96=119$; $S_5=119+600=719$.
  \item Los factoriales $(n+1)!$ son $2,6,24,120,720$: en cada caso $S_n=(n+1)!-1$. Conjetura: $S_n=(n+1)!-1$.
  \item \emph{Caso base:} $S_1=1=2!-1$. \emph{Paso:} si $S_n=(n+1)!-1$, entonces $S_{n+1}=(n+1)!-1+(n+1)(n+1)!=(n+1)!\,(1+n+1)-1=(n+2)!-1$. Por inducción, la conjetura es cierta para todo $n\in\N$. $\blacksquare$
\end{enumerate}
\resultado $S_n=(n+1)!-1$; valores $1,5,23,119,719$.
\comprobacion la identidad $k\cdot k!=(k+1)!-k!$ hace la suma telescópica: $S_n=(n+1)!-1!$. Valores comprobados con Python.
\end{solucion}

\begin{solucion}{eje:5-sumas}
\datos $\sum_{k=1}^{n}(3k-1)$; $\sum_{k=0}^{7}2^{k}$.
\planteamiento linealidad del sumatorio y \cref{ej:gauss}; \cref{ej:geometrica} con $r=2$.
\pasos
\begin{enumerate}[label=\arabic*.]
  \item $\sum_{k=1}^{n}(3k-1)=3\sum_{k=1}^{n}k-\sum_{k=1}^{n}1=\frac{3n(n+1)}{2}-n=\frac{3n^{2}+n}{2}=\frac{n(3n+1)}{2}$.
  \item Inducción. \emph{Caso base:} $n=1$: $2=\frac{1\cdot 4}{2}$. \emph{Paso:} si vale para $n$, $\frac{n(3n+1)}{2}+3(n+1)-1=\frac{3n^{2}+n+6n+4}{2}=\frac{3n^{2}+7n+4}{2}=\frac{(n+1)(3n+4)}{2}$, que es la fórmula para $n+1$. $\blacksquare$
  \item $\sum_{k=0}^{7}2^{k}=\frac{2^{8}-1}{2-1}=255$.
  \item Un byte tiene 8 bits; el mayor número representable es $11111111_2=2^{7}+\dots+2^{0}=255$, y hay $256$ valores, de $0$ a $255$.
\end{enumerate}
\resultado (a) $\dfrac{n(3n+1)}{2}$; (b) $255$.
\comprobacion (a) con $n=4$: $2+5+8+11=26=\frac{4\cdot 13}{2}$. (b) $1+2+4+\dots+128=255$; en hexadecimal, \texttt{0xFF} (módulo~12).
\end{solucion}

\begin{solucion}{eje:5-log}
\datos $p_1,\dots,p_n>0$; $\ln(ab)=\ln a+\ln b$ para $a,b>0$.
\planteamiento inducción sobre $n$, usando la definición recursiva del productorio, $\prod_{t=1}^{n+1}p_t=\bigl(\prod_{t=1}^{n}p_t\bigr)p_{n+1}$.
\pasos
\begin{enumerate}[label=\arabic*.]
  \item Sea $P(n)$: «para cualesquiera reales positivos $p_1,\dots,p_n$, $\ln\prod_{t=1}^{n}p_t=\sum_{t=1}^{n}\ln p_t$». \emph{Caso base:} $n=1$, ambos miembros valen $\ln p_1$.
  \item \emph{Paso.} Sea $n\ge 1$, supóngase $P(n)$ y sean $p_1,\dots,p_{n+1}>0$. El número $a=\prod_{t=1}^{n}p_t$ es positivo (producto de positivos), y $b=p_{n+1}>0$.
  \item Por la propiedad del logaritmo y la hipótesis de inducción, $\ln\prod_{t=1}^{n+1}p_t=\ln(ab)=\ln a+\ln b=\sum_{t=1}^{n}\ln p_t+\ln p_{n+1}=\sum_{t=1}^{n+1}\ln p_t$.
  \item Por inducción, $P(n)$ vale para todo $n\in\N$. $\blacksquare$
  \item En el módulo~1, $P(x_{1:T})=\prod_{t}P(x_t\mid x_{1:t-1})$ es un producto de $T$ factores en $(0,1]$, que para $T$ grande es un número diminuto (por ejemplo, $0{,}1^{1000}=10^{-1000}$, que no cabe en coma flotante). Su logaritmo es una suma de $T$ términos moderados ($1000\cdot\ln 0{,}1\approx -2302{,}6$), con la que se trabaja sin problemas numéricos; además, maximizar una suma es más sencillo que maximizar un producto.
\end{enumerate}
\resultado \textbf{el logaritmo de un producto finito es la suma de los logaritmos.}
\comprobacion $\ln(2\cdot 3\cdot 5)=\ln 30\approx 3{,}4012$ y $\ln 2+\ln 3+\ln 5\approx 0{,}6931+1{,}0986+1{,}6094=3{,}4012$.
\end{solucion}

% ---------------------------------------------------------------------
\section{Tema 6}

\begin{solucion}{eje:6-clasificar}
\datos cuatro enunciados.
\planteamiento identificar la forma lógica (\cref{tema:2}) y aplicar las reglas de demostración y de refutación del \cref{tema:6}.
\pasos
\begin{enumerate}[label=(\alph*)]
  \item Universal, $\forall p\,(p\text{ primo}\wedge p>2\Rightarrow p\text{ impar})$. Se demuestra en general; se refutaría con un primo par mayor que~2.
  \item Existencial. Se demuestra con un testigo ($101$); se refutaría probando que ningún entero entre $100$ y $110$ es primo.
  \item Existencia y unicidad. Se demuestra con un testigo ($x=2$) y un argumento de unicidad; se refutaría mostrando que no hay solución o que hay dos.
  \item Universal negativo, $\forall n\in\Z\ (n^{2}\neq 3)$. Se demuestra en general; se refutaría con un entero $n$ con $n^{2}=3$.
\end{enumerate}
\resultado (a) y (d) universales; (b) existencial; (c) de existencia y unicidad.
\comprobacion los cuatro son verdaderos: en (d), $1^{2}<3<2^{2}$ y $\abs{n}\ge 2\Rightarrow n^{2}\ge 4$.
\end{solucion}

\begin{solucion}{eje:6-contraejemplos}
\datos cuatro enunciados universales falsos.
\planteamiento dar un caso concreto que cumpla la hipótesis y no la tesis.
\pasos
\begin{enumerate}[label=(\alph*)]
  \item $x=\tfrac12$: $x^{2}=\tfrac14<\tfrac12$. (También $x=0$ o $x=1$, donde $x^{2}=x$.)
  \item $a=b=1$: $\sqrt2\neq 2$.
  \item $n=2$: $n^{2}=4$ es múltiplo de~4 y $2$ no lo es.
  \item $n=3$: $2^{3}+1=9=3\cdot 3$.
\end{enumerate}
\resultado contraejemplos $x=\tfrac12$; $a=b=1$; $n=2$; $n=3$.
\comprobacion en (d), $2^{n}+1$ es primo para $n=1,2,4$ ($3,5,17$), lo que podría inducir a error, como en el \cref{ej:euler}.
\end{solucion}

\begin{solucion}{eje:6-demostrar-refutar}
\datos cuatro enunciados.
\planteamiento decidir primero con ejemplos y después demostrar o refutar.
\pasos
\begin{enumerate}[label=(\alph*)]
  \item \textbf{Falso}: $\sqrt2$ y $-\sqrt2$ son irracionales (el segundo por el \cref{eje:4-racional-irracional} con $a=-1$) y su suma es $0$.
  \item \textbf{Verdadero}: si $a=m/n$ y $b=p/q$ con $n,q\neq 0$, entonces $ab=mp/(nq)$ con $nq\neq 0$.
  \item \textbf{Falso}: $n^{2}+n+1=(n^{2}+n)+1$, y $n^{2}+n$ es par (\cref{prop:consecutivos}), luego $n^{2}+n+1$ es impar para todo $n$. Se refuta un existencial con una demostración universal.
  \item \textbf{Falso}: $x=-1$ no tiene raíz cuadrada real, pues $y^{2}\ge 0>-1$.
\end{enumerate}
\resultado (a) falso; (b) verdadero; (c) falso; (d) falso.
\comprobacion en (c) no basta probar $n=1,2,3$ ($3,7,13$, impares): hace falta el argumento general.
\end{solucion}

\begin{solucion}{eje:6-testigos}
\datos dos enunciados existenciales.
\planteamiento dar testigos concretos y verificarlos.
\pasos
\begin{enumerate}[label=\arabic*.]
  \item (a) Sean $a=\tfrac32$ y $b=\tfrac12$. Entonces $\floor{a}=1$, $\ceil{b}=1$ y $\floor{a}\ceil{b}=1$, mientras que $ab=\tfrac34$ y $\floor{ab}=0$. Como $1\neq 0$, el par $(a,b)$ es testigo.
  \item (b) Sean $a=9$, $b=6$, $c=4$, $d=2$. Se cumple $9>6>4>2>0$ y $81+36+16+4=137$. $\blacksquare$
\end{enumerate}
\resultado (a) $a=\tfrac32$, $b=\tfrac12$; (b) $(9,6,4,2)$.
\comprobacion una búsqueda exhaustiva con Python muestra que $(9,6,4,2)$ es la única cuaterna estrictamente decreciente de naturales positivos con suma de cuadrados $137$. En (a), con $a=b=\tfrac12$ no sirve: $0\cdot 1=0=\floor{1/4}$.
\end{solucion}

\begin{solucion}{eje:6-irracional-entre}
\datos $r,s\in\Q$, $r<s$.
\planteamiento construir el testigo $x=r+\dfrac{s-r}{\sqrt2}$.
\pasos
\begin{enumerate}[label=\arabic*.]
  \item Sean $r<s$ racionales y $x=r+\dfrac{s-r}{\sqrt2}$.
  \item Como $s-r>0$ y $\sqrt2>1$, se tiene $0<\dfrac{s-r}{\sqrt2}<s-r$, luego $r<x<s$.
  \item $x$ es irracional: $\dfrac{1}{\sqrt2}$ es irracional (si fuera racional, su inverso $\sqrt2$ también lo sería); por el \cref{eje:4-racional-irracional}, $(s-r)\cdot\dfrac{1}{\sqrt2}$ es irracional por ser $s-r$ racional no nulo, y $r+\dfrac{s-r}{\sqrt2}$ es irracional por ser suma de un racional y un irracional.
  \item Por tanto, $x$ es un irracional entre $r$ y $s$. $\blacksquare$
\end{enumerate}
\resultado \textbf{entre dos racionales distintos hay un irracional}, por ejemplo $r+(s-r)/\sqrt2$.
\comprobacion $r=0$, $s=1$: $x=1/\sqrt2\approx 0{,}7071\in(0,1)$. La prueba es constructiva.
\end{solucion}

\begin{solucion}{eje:6-paridad-unica}
\datos $n\in\Z$; \cref{teo:division} con $b=2$.
\planteamiento existencia y unicidad del cociente; después, reducción al absurdo directa para «par e impar».
\pasos
\begin{enumerate}[label=\arabic*.]
  \item \emph{Existencia.} Por el \cref{teo:division} con $a=n$ y $b=2$, existen $k,r\in\Z$ con $n=2k+r$ y $r\in\{0,1\}$; es decir, $n=2k$ o $n=2k+1$.
  \item \emph{Unicidad.} Si $k$ y $k'$ cumplen la condición, entonces $n=2k+r=2k'+r'$ con $r,r'\in\{0,1\}$, y la unicidad de la división euclídea (\cref{prop:division-unicidad}) da $k=k'$ (y $r=r'$).
  \item \emph{Ningún entero es par e impar} (prueba directa). Supóngase, por reducción al absurdo, que $n=2k=2\ell+1$ con $k,\ell\in\Z$. Entonces $2(k-\ell)=1$, luego $k-\ell=\tfrac12$. Pero $k-\ell$ es entero y no hay enteros estrictamente entre $0$ y $1$. Contradicción. $\blacksquare$
\end{enumerate}
\resultado \textbf{$k$ existe y es único; ningún entero es par e impar.}
\comprobacion $n=-7=2(-4)+1$: $k=-4$ (no $k=-3$, que daría $-6+1=-5$).
\end{solucion}

\begin{solucion}{eje:6-terna}
\datos $p$, $p+2$, $p+4$ primos.
\planteamiento existencia con $p=3$; unicidad con los restos módulo~3.
\pasos
\begin{enumerate}[label=\arabic*.]
  \item \emph{Existencia.} $p=3$: $3$, $5$ y $7$ son primos.
  \item \emph{Unicidad.} Sea $p$ un primo tal que $p,p+2,p+4$ son primos. Por la división euclídea, $p\equiv 0$, $1$ o $2\pmod 3$.
  \item Si $p\equiv 1$, entonces $p+2\equiv 0\pmod 3$, y como $p+2>3$, $p+2$ sería compuesto. Si $p\equiv 2$, entonces $p+4\equiv 6\equiv 0\pmod 3$ y $p+4>3$ sería compuesto. Ambos casos son imposibles.
  \item Luego $p\equiv 0\pmod 3$, es decir, $3\divides p$, y como $p$ es primo, $p=3$. $\blacksquare$
\end{enumerate}
\resultado \textbf{$p=3$ es el único} (terna $3,5,7$).
\comprobacion búsqueda con Python hasta $1000$: solo $p=3$. Entre tres números $p,p+2,p+4$ uno siempre es múltiplo de~3, porque sus restos módulo~3 son $p,p+2,p+1$, distintos.
\end{solucion}

\begin{solucion}{eje:6-mersenne}
\datos (a) $n$ primo $\Rightarrow 2^{n}-1$ primo; (b) $2^{n}-1$ primo $\Rightarrow n$ primo.
\planteamiento (a) buscar contraejemplo; (b) contrarrecíproco con la factorización indicada.
\pasos
\begin{enumerate}[label=\arabic*.]
  \item (a) \textbf{Falso}: $n=11$ es primo y $2^{11}-1=2047=23\cdot 89$.
  \item (b) Se demuestra el contrarrecíproco: si $n\ge 2$ no es primo, $2^{n}-1$ no es primo. Sea $n=ab$ con $1<a,b<n$.
  \item Con $x=2^{a}$ y $m=b$: $2^{n}-1=(2^{a})^{b}-1=(2^{a}-1)\bigl(2^{a(b-1)}+\dots+2^{a}+1\bigr)$.
  \item Como $a\ge 2$, $2^{a}-1\ge 3>1$; como $a<n$, $2^{a}-1<2^{n}-1$. Luego $2^{a}-1$ es un divisor de $2^{n}-1$ estrictamente entre $1$ y $2^{n}-1$, y $2^{n}-1$ es compuesto. $\blacksquare$
\end{enumerate}
\resultado (a) \textbf{falso} ($2^{11}-1=23\cdot 89$); (b) \textbf{verdadero}.
\comprobacion $n=6$: $2^{6}-1=63=(2^{2}-1)(2^{4}+2^{2}+1)=3\cdot 21$. La identidad del paso~3 es la serie geométrica del \cref{ej:geometrica} multiplicada por $x-1$.
\end{solucion}

\begin{solucion}{eje:6-promedio}
\datos $x_1+\dots+x_{10}=100$.
\planteamiento reducción al absurdo.
\pasos
\begin{enumerate}[label=\arabic*.]
  \item Supóngase, por reducción al absurdo, que $x_i<10$ para todo $i$.
  \item Sumando las diez desigualdades, $x_1+\dots+x_{10}<100$, en contra de la hipótesis.
  \item Luego existe $i$ con $x_i\ge 10$. $\blacksquare$
  \item La prueba \textbf{no es constructiva}: no indica qué $x_i$ cumple la propiedad.
  \item No existe necesariamente $i$ con $x_i>10$: con $x_1=\dots=x_{10}=10$ ninguno lo cumple.
\end{enumerate}
\resultado \textbf{existe $x_i\ge 10$} (prueba no constructiva); no necesariamente $x_i>10$.
\comprobacion es el principio del promedio: algún valor es mayor o igual que la media, $100/10=10$.
\end{solucion}

\begin{solucion}{eje:6-critica-1}
\datos la prueba propuesta.
\planteamiento revisar la redacción con las pautas del \cref{tema:6} y la lógica de cada paso.
\pasos
\begin{enumerate}[label=\arabic*.]
  \item \emph{Razonamiento.} Usa la misma $k$ para $x$ e $y$, lo que supone $y=x+1$: solo trata pares consecutivos (\cref{obs:3-errores}(b)).
  \item \emph{Redacción.} No presenta $x$ e $y$ como enteros arbitrarios, y concluye «impar» sin escribir $x+y$ en la forma $2m+1$ con $m$ entero.
  \item El teorema es \textbf{verdadero}. Prueba correcta: «Sean $x$ un entero par e $y$ un entero impar. Existen $k,\ell\in\Z$ con $x=2k$ e $y=2\ell+1$. Entonces $x+y=2(k+\ell)+1$, y como $k+\ell\in\Z$, $x+y$ es impar. $\blacksquare$»
\end{enumerate}
\resultado error de \textbf{variable reutilizada} y conclusión sin justificar; teorema cierto, prueba en el paso~3.
\comprobacion $x=4$, $y=7$: $k=2$, $\ell=3$, $x+y=11=2\cdot 5+1$. La prueba original no cubre este caso, pues $7\neq 4+1$.
\end{solucion}

\begin{solucion}{eje:6-critica-2}
\datos la prueba propuesta.
\planteamiento identificar qué condicional se ha probado.
\pasos
\begin{enumerate}[label=\arabic*.]
  \item El enunciado es $P\Rightarrow Q$ con $P$: $n\equiv 0\pmod 2$ y $Q$: $n\equiv 0\pmod 4$. Su contrarrecíproco es «si $n\not\equiv 0\pmod 4$, entonces $n\not\equiv 0\pmod 2$».
  \item La prueba supone $Q$ y deduce $P$: demuestra el \textbf{recíproco} $Q\Rightarrow P$, que es cierto pero no equivale al enunciado.
  \item El enunciado es \textbf{falso}: $n=2$ cumple $2\equiv 0\pmod 2$ y $2\not\equiv 0\pmod 4$.
\end{enumerate}
\resultado se ha demostrado el recíproco; el enunciado es falso (contraejemplo $n=2$).
\comprobacion es el error de la \cref{obs:3-errores}(d); un contraejemplo al enunciado basta para saber que ninguna prueba correcta puede existir.
\end{solucion}

\begin{solucion}{eje:6-uno-dos}
\datos la cadena de igualdades.
\planteamiento revisar qué operación no es válida.
\pasos
\begin{enumerate}[label=\arabic*.]
  \item Las igualdades $a^{2}=ab$, $a^{2}-b^{2}=ab-b^{2}$ y $(a+b)(a-b)=b(a-b)$ son correctas.
  \item El paso «simplificando» divide ambos miembros por $a-b$, que vale $0$ porque $a=b$. De $X\cdot 0=Y\cdot 0$ no se sigue $X=Y$.
  \item (Si además $b=0$, el último paso divide también por cero.)
\end{enumerate}
\resultado el error es \textbf{dividir por $a-b=0$}.
\comprobacion con $a=b=1$: $(1+1)\cdot 0=1\cdot 0$ es cierto ($0=0$), pero $2\neq 1$.
\end{solucion}

\begin{solucion}{eje:6-redactar}
\datos $a,b\in\N$, $a\divides b$, $b\divides a$.
\planteamiento redacción directa con frases completas y justificación de cada paso.
\pasos
\begin{enumerate}[label=\arabic*.]
  \item Sean $a,b\in\N$ con $a\divides b$ y $b\divides a$. Por la \cref{def:divide}, existen enteros $k$ y $\ell$ con $b=ak$ y $a=b\ell$.
  \item Sustituyendo la primera igualdad en la segunda, $a=ak\ell$. Como $a\ge 1$, en particular $a\neq 0$, y puede dividirse por $a$: $k\ell=1$.
  \item Como $a,b>0$ y $b=ak$, el entero $k$ es positivo, y por la misma razón $\ell>0$. Dos enteros positivos con producto~1 son ambos iguales a~1: si fuera $k\ge 2$, sería $k\ell\ge 2\ell\ge 2$. Así $k=\ell=1$.
  \item Por tanto, $b=a\cdot 1=a$. $\blacksquare$
\end{enumerate}
\resultado \textbf{$a=b$}, con la demostración redactada en los pasos 1–4.
\comprobacion la hipótesis de positividad es necesaria: en $\Z$, $a=2$ y $b=-2$ se dividen mutuamente y son distintos (ahí $k=\ell=-1$). Compárese con el borrador del enunciado: las flechas se han sustituido por frases que explican \emph{por qué} cada paso es válido.
\end{solucion}

\begin{solucion}{eje:6-unica-raiz}
\datos $a\neq 0$; $ax^{2}+bx+c=0$.
\planteamiento completar el cuadrado: $ax^{2}+bx+c=a\Bigl[\bigl(x+\tfrac{b}{2a}\bigr)^{2}-\tfrac{\Delta}{4a^{2}}\Bigr]$ con $\Delta=b^{2}-4ac$.
\pasos
\begin{enumerate}[label=\arabic*.]
  \item Se comprueba la identidad del planteamiento desarrollando: $a\bigl(x^{2}+\tfrac{b}{a}x+\tfrac{b^{2}}{4a^{2}}\bigr)-\tfrac{b^{2}-4ac}{4a}=ax^{2}+bx+c$. Como $a\neq 0$, $x$ es solución si y solo si $\bigl(x+\tfrac{b}{2a}\bigr)^{2}=\tfrac{\Delta}{4a^{2}}$.
  \item ($\Leftarrow$) Si $\Delta=0$, la condición es $\bigl(x+\tfrac{b}{2a}\bigr)^{2}=0$, que equivale a $x=-\tfrac{b}{2a}$. Hay solución y es única.
  \item ($\Rightarrow$) Por contrarrecíproco: si $\Delta\neq 0$, no hay solución única. Si $\Delta<0$, el miembro derecho es negativo y el izquierdo no: no hay solución. Si $\Delta>0$, sea $s=\sqrt{\Delta}/(2\abs{a})>0$; entonces $x=-\tfrac{b}{2a}+s$ y $x=-\tfrac{b}{2a}-s$ son dos soluciones distintas.
  \item En ambos casos falla la existencia o la unicidad. Por tanto, hay solución única si y solo si $\Delta=0$, y entonces es $x=-b/(2a)$. $\blacksquare$
\end{enumerate}
\resultado \textbf{solución única $\iff b^{2}=4ac$}, y es $x=-b/(2a)$.
\comprobacion $x^{2}-6x+9=(x-3)^{2}$: $\Delta=36-36=0$ y $x=6/2=3$. $x^{2}-1$: $\Delta=4>0$, soluciones $\pm1$. $x^{2}+1$: $\Delta=-4$, ninguna.
\end{solucion}
